\chaptermark{Publication exceptionnelle du continu}
\label{chap:83-20}

\input{ampliaciones_sucesoras_20260824/parte_i_ii/tex/bloques/salto_espectral_4_2_6_54}

% Unidad U018. Paley, extension cuadratica finita y codigo ternario.
% Redaccion autonoma desde la autoridad material 1063.

\section{Forme de conférence de Paley, extension quadratique finie et code ternaire autodual}
\label{p12-83-c20-s1:chap:paley-extension-cuadratica-codigo-ternario}

La tétrade \(B_K\) construite dans l'unité précédente ne possède pas encore
d'étoile de Witt : pour la définir, il faut d'abord construire le système
d'incidence qui contient ses hexades. Cette unité réalise cette étape à partir de la
jauge projective des six unités modulo neuf, obtient une racine
quadratique interne de \(-I_6\) et construit le code ternaire en graphe de
longueur douze. Aucune identification ne se déduit d'une simple égalité de
cardinaux.

\subsection{Jauge projective des \(6\) unités modulo \(9\)}

Soit
\[
 \mathcal U_9=(\mathbb Z/9\mathbb Z)^\times
 =\{1,2,4,8,7,5\},
\]
ordonné par les puissances successives de \(2\) modulo \(9\). Nous fixons la
bijection de coordonnées
\begin{equation}
 \theta:\mathcal U_9\longrightarrow\mathbb P^1(\mathbb F_5),
 \qquad
 \begin{array}{c|cccccc}
 u&1&2&4&8&7&5\\ \hline
 \theta(u)&\infty&0&1&2&3&4
 \end{array}.
 \label{p12-83-c20-s1:eq:calibre-unidades-paley}
\end{equation}
L'application \(\theta\) est une jauge de coordonnées. Il n'est pas affirmé
qu'elle soit un homomorphisme entre la multiplication de \(\mathcal U_9\) et une
opération sur \(\mathbb P^1(\mathbb F_5)\).

Soit
\[
 \chi:\mathbb F_5\longrightarrow\{-1,0,1\}
\]
le caractère quadratique :
\[
 \chi(0)=0,
 \qquad
 \chi(1)=\chi(4)=1,
 \qquad
 \chi(2)=\chi(3)=-1.
\]
Dans l'ordre \((\infty,0,1,2,3,4)\), nous définissons
\begin{align}
 (C_P)_{\infty,x}&=(C_P)_{x,\infty}=1,
 &&x\in\mathbb F_5,\label{p12-83-c20-s1:eq:paley-fila-infinito}\\
 (C_P)_{x,y}&=\chi(x-y),
 &&x,y\in\mathbb F_5, x\neq y,\label{p12-83-c20-s1:eq:paley-entradas-finitas}\\
 (C_P)_{x,x}&=0.
 \label{p12-83-c20-s1:eq:paley-diagonal}
\end{align}
Ainsi,
\begin{equation}
 C_P=
 \begin{pmatrix}
 0& 1& 1& 1& 1& 1\\
 1& 0& 1&-1&-1& 1\\
 1& 1& 0& 1&-1&-1\\
 1&-1& 1& 0& 1&-1\\
 1&-1&-1& 1& 0& 1\\
 1& 1&-1&-1& 1& 0
 \end{pmatrix}.
 \label{p12-83-c20-s1:eq:matriz-conferencia-paley-autonoma}
\end{equation}

\begin{lemma}[Somme de corrélation quadratique]
\label{p12-83-c20-s1:lem:correlacion-cuadratica-f5}
Pour \(a\neq b\) dans \(\mathbb F_5\),
\begin{equation}
 \sum_{t\in\mathbb F_5}
 \chi(t-a)\chi(t-b)=-1.
 \label{p12-83-c20-s1:eq:correlacion-cuadratica-f5}
\end{equation}
\end{lemma}

\begin{proof}
La substitution \(u=(t-a)/(b-a)\) est une bijection de \(\mathbb F_5\).
Puisque \(\chi((b-a)^2)=1\), la somme se réduit à
\[
 \sum_{u\in\mathbb F_5}\chi(u(u-1)).
\]
L'évaluation de \(u=0,1,2,3,4\) donne respectivement
\(0,0,1,-1,-1\), dont la somme est \(-1\).
\end{proof}

\begin{proposition}[Identité de conférence]
\label{p12-83-c20-s1:prop:paley-identidad-conferencia-autonoma}
La matrice \(C_P\) est symétrique et satisfait
\begin{equation}
 \boxed{C_PC_P^{\mathsf T}=5I_6.}
 \label{p12-83-c20-s1:eq:paley-identidad-conferencia-autonoma}
\end{equation}
\end{proposition}

\begin{proof}
Chaque ligne contient cinq entrées de module un et une entrée nulle ; le carré de sa
norme est cinq. Pour deux lignes finies distinctes, la contribution
de la coordonnée \(\infty\) est \(+1\), tandis que
\cref{p12-83-c20-s1:lem:correlacion-cuadratica-f5} donne \(-1\) pour les cinq coordonnées
finies. Leur produit scalaire est nul. La somme d'un caractère non trivial
sur \(\mathbb F_5\) est nulle, de sorte que la ligne correspondant à
\(\infty\) est également orthogonale à chaque ligne finie.
\end{proof}

\subsection{Réduction ternaire et racine quadratique interne de \texorpdfstring{\(-I_6\)}{-I₆}}

Nous réduisons \(C_P\) modulo trois, avec \(-1\equiv2\pmod3\) :
\begin{equation}
 A_W=
 \begin{pmatrix}
 0&1&1&1&1&1\\
 1&0&1&2&2&1\\
 1&1&0&1&2&2\\
 1&2&1&0&1&2\\
 1&2&2&1&0&1\\
 1&1&2&2&1&0
 \end{pmatrix}
 \in M_6(\mathbb F_3).
 \label{p12-83-c20-s1:eq:matriz-paley-witt-ternaria}
\end{equation}

\begin{theorem}[Racine quadratique interne de \(-I_6\)]
\label{p12-83-c20-s1:thm:aw-cuadrado-menos-identidad-autonomo}
On a
\begin{equation}
 \boxed{A_W^2=-I_6.}
 \label{p12-83-c20-s1:eq:aw-cuadrado-menos-identidad-autonomo}
\end{equation}
En particulier,
\[
 A_W^{-1}=-A_W.
\]
\end{theorem}

\begin{proof}
La \cref{p12-83-c20-s1:prop:paley-identidad-conferencia-autonoma} donne
\[
 A_WA_W^{\mathsf T}
 \equiv5I_6
 \equiv2I_6
 =-I_6
 \pmod3.
\]
Puisque \(A_W=A_W^{\mathsf T}\), on obtient \(A_W^2=-I_6\).
\end{proof}

L'identité précédente est l'espace de résidence algébrique de la racine de \(-1\)
dans ce corridor. Elle ne provient pas du facteur \(1/4\) de l'inversion de Hadamard :
celui-ci est un scalaire rationnel imposé par \(H_4^2=4I_4\) ; celle-là est une
structure quadratique en caractéristique trois.

\subsection{Extension quadratique à \texorpdfstring{\(\mathbb F_9\)}{F₉}}

On définit
\begin{equation}
 \mathbb F_9=\mathbb F_3[\iota]/(\iota^2+1),
 \qquad \iota^2=-1.
 \label{p12-83-c20-s1:eq:definicion-f9-autonoma}
\end{equation}
Le polynôme \(X^2+1\) n'a pas de racines dans \(\mathbb F_3\), de sorte que le
quotient est un corps à neuf éléments. Pour
\(V=\mathbb F_3^6\), soit
\[
 V_9=V\otimes_{\mathbb F_3}\mathbb F_9.
\]
Les opérateurs
\begin{equation}
 P_\iota=\frac12(I-\iota A_W),
 \qquad
 P_{-\iota}=\frac12(I+\iota A_W)
 \label{p12-83-c20-s1:eq:proyectores-f9-paley}
\end{equation}
sont bien définis puisque \(2\) est inversible dans \(\mathbb F_3\).

\begin{proposition}[Projecteurs conjugués]
\label{p12-83-c20-s1:prop:proyectores-conjugados-paley}
On a
\[
 P_\iota^2=P_\iota,
 \qquad
 P_{-\iota}^2=P_{-\iota},
 \qquad
 P_\iota P_{-\iota}=0,
 \qquad
 P_\iota+P_{-\iota}=I.
\]
\end{proposition}

\begin{proof}
En utilisant \(\iota^2=-1\) et \(A_W^2=-I\),
\[
 (\iota A_W)^2=I.
\]
Les quatre identités sont alors les identités élémentaires des
projecteurs \((I\mp\iota A_W)/2\) d'une involution.
\end{proof}

\begin{theorem}[Décomposition spectrale conjuguée]
\label{p12-83-c20-s1:thm:descomposicion-espectral-f9-paley}
Soit
\[
 E_{\pm\iota}=\ker(A_W\mp\iota I).
\]
Alors
\begin{equation}
 V_9=E_\iota\oplus E_{-\iota},
 \qquad
 \dim_{\mathbb F_9}E_\iota
 =\dim_{\mathbb F_9}E_{-\iota}=3.
 \label{p12-83-c20-s1:eq:descomposicion-espectral-f9-paley}
\end{equation}
\end{theorem}

\begin{proof}
Le polynôme minimal de \(A_W\) divise \(X^2+1\), qui se scinde sur
\(\mathbb F_9\) avec des racines simples \(\pm\iota\). Les projecteurs de
\cref{p12-83-c20-s1:prop:proyectores-conjugados-paley} produisent la somme directe. Puisque
\(A_W\) a ses coefficients dans \(\mathbb F_3\), l'automorphisme de
Frobenius échange les deux espaces propres ; leurs dimensions sont égales et
leur somme vaut six.
\end{proof}

Cette construction dote \(\mathbb F_3^6\) d'une structure de module de rang
trois sur \(\mathbb F_9\), au moyen de
\begin{equation}
 (a+b\iota)\cdot v:=av+b(vA_W).
 \label{p12-83-c20-s1:eq:accion-f9-sobre-f3-seis}
\end{equation}
L'égalité \(3^6=9^3\) est ainsi accompagnée d'une action explicite ; elle n'est pas
utilisée comme argument cardinal.

\subsection{Code ternaire en graphe de longueur \(12\)}

On définit
\begin{equation}
 c:\mathbb F_3^6\longrightarrow\mathbb F_3^{12},
 \qquad
 c(q)=(q,qA_W),
 \label{p12-83-c20-s1:eq:aplicacion-codigo-grafico-witt}
\end{equation}
et
\begin{equation}
 C_W:=\operatorname{im}c.
 \label{p12-83-c20-s1:eq:definicion-codigo-witt-ternario}
\end{equation}
Sa matrice génératrice est
\[
 G_W=(I_6\mid A_W).
\]

\begin{theorem}[Injectivité et autodualité]
\label{p12-83-c20-s1:thm:cw-inyectivo-autodual}
L'application \(c\) est injective, \(\dim C_W=6\) et
\begin{equation}
 \boxed{C_W=C_W^\perp.}
 \label{p12-83-c20-s1:eq:cw-autodual}
\end{equation}
\end{theorem}

\begin{proof}
La première moitié de \(c(q)\) est \(q\), de sorte que
\(\operatorname{pr}_1\circ c=\operatorname{id}\) ; cela prouve
l'injectivité et la dimension. En outre,
\[
 G_WG_W^{\mathsf T}
 =I_6+A_WA_W^{\mathsf T}
 =I_6+2I_6=0
\]
dans \(\mathbb F_3\). Par conséquent, \(C_W\subseteq C_W^\perp\). Les deux espaces
sont de dimension six, donc ils sont égaux.
\end{proof}

Pour \(r=0,\ldots,12\), soit
\[
 N_r:=\#\{q\in\mathbb F_3^6:
            \operatorname{wt}(q,qA_W)=r\}.
\]
L'énumération des \(3^6=729\) entrées produit
\begin{equation}
 N_0=1,
 \qquad
 N_6=264,
 \qquad
 N_9=440,
 \qquad
 N_{12}=24,
 \label{p12-83-c20-s1:eq:censo-pesos-cw}
\end{equation}
et \(N_r=0\) pour les autres poids.

\begin{theorem}[Énumérateur et distance minimale]
\label{p12-83-c20-s1:thm:enumerador-cw-autonomo}
Le code \(C_W\) possède les paramètres
\[
 \boxed{C_W=[12,6,6]_3}
\]
et l'énumérateur des poids
\begin{equation}
 \boxed{W_{C_W}(y)=1+264y^6+440y^9+24y^{12}.}
 \label{p12-83-c20-s1:eq:enumerador-pesos-cw}
\end{equation}
\end{theorem}

\begin{proof}
La dimension et l'autodualité ont déjà été démontrées. Le recensement
\eqref{p12-83-c20-s1:eq:censo-pesos-cw} parcourt les \(729\) mots sans échantillonnage. Son
plus petit poids non nul est six, qui est la distance minimale d'un code linéaire.
La somme \(1+264+440+24=729\) contrôle qu'aucun mot n'a été omis.
\end{proof}

\subsection{Nature de la structure obtenue}

Trois objets distincts et trois flèches typées ont été construits :
\[
 \mathcal U_9
 \xrightarrow{\theta}\mathbb P^1(\mathbb F_5)
 \xrightarrow{\chi}C_P
 \xrightarrow{\bmod3}A_W,
\]
\[
 A_W^2=-I_6
 \quad\Longrightarrow\quad
 \mathbb F_3^6\cong\mathbb F_9^3
 \quad\Longrightarrow\quad
 C_W=\{(q,qA_W):q\in\mathbb F_3^6\}.
\]
La première flèche est une jauge, la deuxième utilise le caractère quadratique et la
troisième est une réduction des coefficients. La structure sur
\(\mathbb F_9\) est finie ; elle ne s'identifie pas à une structure complexe
réelle ou analytique.

\subsection{Obstructions du corridor Paley–\(C_W\) : bijection, \(A_W^2=-I_6\), autodualité et recensement de \(729\) mots}

\begin{criterion}[Réfutation du corridor Paley--\(C_W\)]
La construction est réfutée si l'un des faits
suivants se produit :
\begin{enumerate}
 \item \(\theta\) n'est pas bijective ou est utilisée comme homomorphisme sans preuve ;
 \item la somme \eqref{p12-83-c20-s1:eq:correlacion-cuadratica-f5} diffère de \(-1\) ;
 \item \(C_PC_P^{\mathsf T}\neq5I_6\);
 \item \(A_W^2\neq-I_6\);
 \item les projecteurs \eqref{p12-83-c20-s1:eq:proyectores-f9-paley} ne sont pas
       complémentaires ;
 \item l'application \(q\mapsto(q,qA_W)\) perd l'injectivité ou
       l'autodualité ;
 \item la somme du recensement ne vaut pas \(729\), diffère de
       \eqref{p12-83-c20-s1:eq:censo-pesos-cw}, ou présente un mot de poids compris entre un
       et cinq ;
 \item on infère \(\mathbb F_3^6\cong\mathbb F_9^3\) uniquement de
       \(3^6=9^3\), en omettant l'action \eqref{p12-83-c20-s1:eq:accion-f9-sobre-f3-seis}.
\end{enumerate}
\end{criterion}


% Unidad U019. Sistema de Witt, campo de estrellas y grupo M12.
% Redaccion autonoma desde la autoridad material 1063.

\section{Système de Steiner ternaire, involutions d'étoile de Witt et génération de \texorpdfstring{\(M_{12}\)}{M12}}
\label{p12-83-c20-s2:chap:sistema-witt-estrellas-m12}

L'unité précédente a construit le code autodual
\(C_W=[12,6,6]_3\). On projectivise maintenant ses mots de poids six,
on prouve la propriété de Steiner, et c'est alors seulement que l'on définit l'étoile
associée à la tétrade \(B_K\). Cet ordre empêche de transformer une figure
d'incidence en un objet ajouté par ressemblance visuelle.

\subsection{Supports projectifs de poids \(6\)}

Soit
\begin{equation}
 \mathcal H_W
 :=\{\operatorname{supp}(c):c\in C_W,
                              \operatorname{wt}(c)=6\}.
 \label{p12-83-c20-s2:eq:familia-soportes-peso-seis}
\end{equation}
Chaque support apparaît pour la paire de mots \(\{c,-c\}\). En effet, si deux
mots de poids six ayant le même support n'étaient pas proportionnels, on pourrait
multiplier l'un d'eux par un scalaire pour annuler une coordonnée en les soustrayant ; on obtiendrait
un mot non nul de \(C_W\) de poids inférieur à six, ce qui contredit la distance
minimale. Puisqu'il n'existe pas non plus de mot non nul égal à son opposé en
caractéristique trois, chaque orbite projective possède exactement deux éléments.
Par conséquent,
\begin{equation}
 |\mathcal H_W|=\frac{264}{2}=132.
 \label{p12-83-c20-s2:eq:numero-hexadas-cw}
\end{equation}

Pour \(T\in\binom{[12]}5\), nous définissons
\begin{equation}
 \nu(T):=\#\{H\in\mathcal H_W:T\subseteq H\}.
 \label{p12-83-c20-s2:eq:incidencia-cinco-hexada}
\end{equation}
L'énumération exacte des
\[
 \binom{12}{5}=792
\]
sous-ensembles de cinq points produit
\begin{equation}
 \nu(T)=1
 \qquad
 \text{pour tout }T\in\binom{[12]}5.
 \label{p12-83-c20-s2:eq:unicidad-cinco-en-hexada}
\end{equation}

\begin{theorem}[Système de Witt ternaire]
\label{p12-83-c20-s2:thm:w12-autonomo}
Le système d'incidence
\[
 W_{12}:=([12],\mathcal H_W)
\]
est le système de Steiner
\begin{equation}
 \boxed{W_{12}=S(5,6,12).}
 \label{p12-83-c20-s2:eq:w12-steiner-autonomo}
\end{equation}
\end{theorem}

\begin{proof}
L'égalité \eqref{p12-83-c20-s2:eq:unicidad-cinco-en-hexada} est précisément la
propriété de Steiner. Le nombre de blocs se retrouve également par le double
comptage des couples \((T,H)\) avec \(T\subset H\) :
\[
 |\mathcal H_W|\binom65=\binom{12}{5},
 \qquad
 |\mathcal H_W|=\frac{792}{6}=132.
\]
\end{proof}

\subsection{Résidu d'une tétrade et décomposition en \(4\) composantes incidentes}

Soit \(B\in\binom{[12]}4\). Dans un système \(S(5,6,12)\), le nombre
d'hexades qui contiennent \(B\) est
\begin{equation}
 \lambda_4
 =\frac{\binom{12-4}{5-4}}{\binom{6-4}{5-4}}
 =\frac82=4.
 \label{p12-83-c20-s2:eq:lambda-cuatro-w12}
\end{equation}
Nous écrivons ces quatre hexades sous la forme
\[
 H_i=B\sqcup P_i,
 \qquad |P_i|=2,
 \qquad i=1,\ldots,4.
\]

\begin{lemma}[Partition résiduelle d'une tétrade]
\label{p12-83-c20-s2:lem:particion-residual-tetrada-witt}
Les quatre paires \(P_1,\ldots,P_4\) sont disjointes et
\begin{equation}
 [12]\setminus B
 =P_1\sqcup P_2\sqcup P_3\sqcup P_4.
 \label{p12-83-c20-s2:eq:particion-residual-tetrada-witt}
\end{equation}
\end{lemma}

\begin{proof}
Si deux paires résiduelles partageaient un point \(p\), les hexades
correspondantes contiendraient le même sous-ensemble \(B\cup\{p\}\) de cinq points,
en contradiction avec l'unicité de Steiner. Les quatre paires sont donc
disjointes ; elles contiennent huit points et épuisent le complémentaire de \(B\).
\end{proof}

\begin{definition}[Étoile et involution d'étoile de Witt]
\label{p12-83-c20-s2:def:estrella-witt-involucion}
L'étoile centrée en \(B\) est
\begin{equation}
 \mathscr S_B
 :=\{H\in\mathcal H_W:B\subseteq H\}
 =\{B\sqcup P_1,\ldots,B\sqcup P_4\}.
 \label{p12-83-c20-s2:eq:estrella-witt-general}
\end{equation}
Si \(P_i=\{p_i^-,p_i^+\}\), son involution est
\begin{equation}
 \sigma_B:=\prod_{i=1}^{4}(p_i^-\ p_i^+).
 \label{p12-83-c20-s2:eq:involucion-estrella-witt-general}
\end{equation}
La permutation \(\sigma_B\) fixe point par point \(B\) et échange les
extrémités de chaque pétale résiduel.
\end{definition}

Le terme mathématique employé désormais est \emph{étoile de Witt} ou
\emph{involution d'étoile de Witt}. Il ne s'agit pas d'une étoile
euclidienne.

\subsection{Étoile sélectionnée par la tétrade du sceau}

L'unité dodécaphasique a produit, sans utiliser le plan de Steiner,
\[
 B_K=\{4,6,9,10\}.
\]
Maintenant que \(W_{12}\) est construit, son étoile est bien définie. Les
quatre hexades incidentes sont
\begin{align}
 H_{13}&=B_K\sqcup\{1,3\}
       =\{1,3,4,6,9,10\},
 \label{p12-83-c20-s2:eq:hexada-petalo-13}\\
 H_{2,11}&=B_K\sqcup\{2,11\}
       =\{2,4,6,9,10,11\},
 \label{p12-83-c20-s2:eq:hexada-petalo-211}\\
 H_{5,7}&=B_K\sqcup\{5,7\}
       =\{4,5,6,7,9,10\}
       =H_\alpha^-,
 \label{p12-83-c20-s2:eq:hexada-petalo-57}\\
 H_{8,12}&=B_K\sqcup\{8,12\}
       =\{4,6,8,9,10,12\}.
 \label{p12-83-c20-s2:eq:hexada-petalo-812}
\end{align}
En conséquence,
\begin{equation}
 \boxed{
 \sigma_{B_K}=(1\ 3)(2\ 11)(5\ 7)(8\ 12).}
 \label{p12-83-c20-s2:eq:involucion-estrella-witt-bk}
\end{equation}

\begin{figure}[htbp]
\centering
\begin{tikzpicture}[
  >=Stealth,
  font=\small,
  core/.style={draw,double,rounded corners=2mm,align=center,
    minimum width=4.0cm,minimum height=1.3cm,fill=black!4},
  petal/.style={draw,rounded corners=2mm,align=center,
    minimum width=3.55cm,minimum height=1.08cm,fill=black!2},
  negative/.style={petal,line width=1.1pt,double,fill=black!10},
  arr/.style={->,thick}
]
\node[core] (B) at (0,0)
 {\(\displaystyle B_K=\{4,6,9,10\}\)\\
  points fixés par \(\sigma_{B_K}\)};
\node[petal] (P13) at (0,3.0)
 {\(P_1=\{1,3\}\), \((1\ 3)\)\\
  \(H_{13}=B_K\sqcup P_1\)};
\node[petal] (P211) at (5.25,0)
 {\(P_2=\{2,11\}\), \((2\ 11)\)\\
  \(H_{2,11}=B_K\sqcup P_2\)};
\node[negative] (P57) at (0,-3.0)
 {\(P_3=\{5,7\}\), \((5\ 7)\)\\
  \(H_{5,7}=H_\alpha^-\)};
\node[petal] (P812) at (-5.25,0)
 {\(P_4=\{8,12\}\), \((8\ 12)\)\\
  \(H_{8,12}=B_K\sqcup P_4\)};
\draw[arr] (B)--(P13);
\draw[arr] (B)--(P211);
\draw[arr] (B)--(P57);
\draw[arr] (B)--(P812);
\node[align=center,font=\footnotesize] at (0,-4.35)
 {\(\mathscr S_{B_K}
   =\{H_{13},H_{2,11},H_{5,7},H_{8,12}\}\).\\
  Le diagramme représente l'incidence, non la distance géométrique.};
\end{tikzpicture}
\caption{Étoile de Witt centrée en \(B_K\). Le centre est la tétrade
fixe ; les quatre pétales sont les paires résiduelles qui déterminent
l'involution.}
\label{p12-83-c20-s2:fig:estrella-witt-bk-autonoma}
\end{figure}

\subsection{Champ global des 495 étoiles}

Il existe une étoile pour chaque tétrade, de sorte que
\begin{equation}
 \#\{\mathscr S_B:B\in\tbinom{[12]}4\}
 =\binom{12}{4}=495.
 \label{p12-83-c20-s2:eq:numero-estrellas-witt}
\end{equation}
Le recensement exact vérifie
\begin{equation}
 \sigma_B(\mathcal H_W)=\mathcal H_W
 \qquad
 \text{pour toute }B\in\binom{[12]}4.
 \label{p12-83-c20-s2:eq:estrellas-conservan-hexadas}
\end{equation}
On définit
\begin{equation}
 \Gamma_\star
 :=\langle\sigma_B:B\in\tbinom{[12]}4\rangle
 \leq\operatorname{Aut}(W_{12}).
 \label{p12-83-c20-s2:eq:grupo-generado-estrellas-witt}
\end{equation}

Un ensemble générateur de six tétrades est
\[
 1234,
 \quad1235,
 \quad1236,
 \quad1245,
 \quad1345,
 \quad2345.
\]
Les ordres des sous-groupes obtenus en ajoutant successivement leurs
involutions sont
\begin{equation}
 1, 2, 6, 18, 360, 7920, 95040.
 \label{p12-83-c20-s2:eq:cadena-ordenes-generadores-m12}
\end{equation}

\begin{theorem}[Génération du groupe de Mathieu \(M_{12}\)]
\label{p12-83-c20-s2:thm:495-estrellas-m12-autonomo}
On a
\begin{equation}
 \boxed{
 \langle\sigma_B:B\in\tbinom{[12]}4\rangle
 =\operatorname{Aut}(W_{12})
 \cong M_{12}.}
 \label{p12-83-c20-s2:eq:estrellas-generan-m12}
\end{equation}
\end{theorem}

\begin{proof}
Chaque \(\sigma_B\) conserve les \(132\) hexades, donc
\(\Gamma_\star\leq\operatorname{Aut}(W_{12})\). Le recensement de Schreier des
six générateurs indiqués produit l'ordre \(95040\). Le groupe
d'automorphismes de \(S(5,6,12)\) est \(M_{12}\) et possède ce même ordre ;
l'inclusion est donc une égalité.
\end{proof}

Une étoile isolée engendre uniquement
\[
 \langle\sigma_B\rangle\cong C_2.
\]
La conclusion concernant \(M_{12}\) appartient au champ global des
\(495\) involutions, non à \(\sigma_{B_K}\) considérée isolément.

\subsection{Action sur le module d'augmentation}

Soit
\[
 V_{11}=\mathbf1^\perp\subset\mathbb R^{12}.
\]
L'involution \(\sigma_B\) fixe les quatre coordonnées de \(B\) et agit
par quatre transpositions sur son complémentaire. Dans \(\mathbb R^{12}\),
elle possède huit directions de valeur propre \(+1\) et quatre de valeur propre \(-1\).
La direction constante appartient au premier espace propre. Par conséquent,
\begin{equation}
 \operatorname{Spec}(\sigma_B|_{V_{11}})
 =\{1^{[7]},(-1)^{[4]}\},
 \qquad
 \frac1{11}\operatorname{Tr}(\sigma_B|_{V_{11}})=\frac3{11}.
 \label{p12-83-c20-s2:eq:espectro-involucion-estrella-aumentacion}
\end{equation}
C'est une involution signée ; ce n'est pas un projecteur positif de rang trois.

\subsection{Obstructions de \(W_{12}\), du champ de \(495\) involutions d'étoile et de la génération de \(M_{12}\)}

\begin{criterion}[Réfutation de \(W_{12}\), de ses étoiles et de \(M_{12}\)]
La construction est réfutée si l'un des faits
suivants se produit :
\begin{enumerate}
 \item la projectivisation des \(264\) mots de poids six ne produit pas
       \(132\) supports ;
 \item il existe un sous-ensemble de cinq points contenu dans zéro hexade ou dans plus
       d'une hexade ;
 \item une tétrade n'appartient pas exactement à quatre hexades ;
 \item les quatre paires résiduelles ne partitionnent pas le complémentaire de la
       tétrade ;
 \item les pétales de \(B_K\) ne sont pas
       \(\{1,3\},\{2,11\},\{5,7\},\{8,12\}\) ;
 \item l'une des \(495\) involutions ne conserve pas
       \(\mathcal H_W\) ;
 \item le groupe engendré n'a pas pour ordre \(95040\) ;
 \item on attribue \(M_{12}\) à une unique étoile au lieu du champ
       global des involutions ;
 \item l'étoile est introduite avant de construire \(W_{12}\).
\end{enumerate}
\end{criterion}


% Unidad U020. Prolongacion binaria W12--W24 y M24.
% Redaccion autonoma desde la autoridad material 1063.

\section{Extension binaire relative à une dodécade, système \texorpdfstring{\(W_{24}\)}{W24} et groupe \texorpdfstring{\(M_{24}\)}{M24}}
\label{p12-83-c20-s3:chap:dodecada-w24-m24}

Le système \(W_{12}\) construit sur douze phases admet une prolongation
binaire lorsqu'on déclare un code de Golay étendu, une dodécade et un
isomorphisme d'incidence. Ces données sont des entrées relatives de cette branche :
elles ne se déduisent pas uniquement de \(C_W\) et n'interviennent pas dans la branche de réseaux
ternaire qui conduira au réseau de Leech.

\subsection{Code binaire de Golay et octades}

Soit
\[
 \mathcal G_{24}\subset\mathbb F_2^{24}
\]
le code binaire étendu de Golay. Ses poids sont
\begin{equation}
 0, 8, 12, 16, 24.
 \label{p12-83-c20-s3:eq:pesos-golay-binario-extendido}
\end{equation}
Les supports des mots de poids huit forment le système de Steiner
\begin{equation}
 W_{24}=S(5,8,24).
 \label{p12-83-c20-s3:eq:w24-steiner-octadas}
\end{equation}
Ses blocs sont appelés \emph{octades}. Fixons une dodécade \(D\), c'est-à-
dire le support d'un mot de poids douze, et soit \(\mathcal O_{24}\)
la famille des octades.

On définit
\begin{equation}
 \mathcal H(D)
 :=\{O\cap D:O\in\mathcal O_{24},\ |O\cap D|=6\}.
 \label{p12-83-c20-s3:eq:hexadas-residuales-dodecada}
\end{equation}

\begin{theorem}[Plan résiduel d'une dodécade]
\label{p12-83-c20-s3:thm:diseno-residual-dodecada-autonomo}
Le système résiduel est
\begin{equation}
 \boxed{(D,\mathcal H(D))\cong S(5,6,12).}
 \label{p12-83-c20-s3:eq:diseno-residual-dodecada-w12}
\end{equation}
En outre, chaque \(H\in\mathcal H(D)\) possède une unique octade \(O_H\) telle que
\begin{equation}
 O_H\cap D=H.
 \label{p12-83-c20-s3:eq:elevacion-unica-hexada-octada}
\end{equation}
\end{theorem}

\begin{proof}
Soit \(A\subset D\), avec \(|A|=5\). Dans \(W_{24}=S(5,8,24)\), il existe une
unique octade \(O_A\) qui contient \(A\). Si \(j=|O_A\cap D|\), la somme
binaire des mots caractéristiques de \(O_A\) et \(D\) appartient à
\(\mathcal G_{24}\) et a pour poids
\[
 8+12-2j=20-2j.
\]
Puisque \(A\subseteq O_A\cap D\), on a \(5\leq j\leq8\). Les quatre
poids possibles \(10,8,6,4\), combinés à
\eqref{p12-83-c20-s3:eq:pesos-golay-binario-extendido}, imposent \(j=6\). Ainsi, tout
sous-ensemble de cinq points de \(D\) appartient à une unique hexade résiduelle.

Si deux octades avaient la même intersection hexadique avec \(D\), toutes deux
contiendraient les mêmes cinq points de cette hexade, ce qui contredirait
l'unicité dans \(S(5,8,24)\). Cela prouve également l'unicité du
relèvement.
\end{proof}

\subsection{Jauge relative entre le plan HMT et la dodécade}

Pour relier la construction précédente au plan obtenu à partir de
\(C_W\), il faut déclarer explicitement
\begin{equation}
 (D,\ell),
 \qquad
 \ell:W_{12}^{\rm HMT}
 \xrightarrow{\sim}(D,\mathcal H(D)),
 \label{p12-83-c20-s3:eq:dato-relativo-dodecada-alineamiento}
\end{equation}
où \(\ell\) est un isomorphisme d'incidence. Si \(\ell_*\) désigne son
action sur les blocs, nous définissons
\begin{equation}
 \widehat\iota_{D,\ell}
 :=\iota_D\circ\ell_*:
 W_{12}^{\rm HMT}\longrightarrow W_{24},
 \qquad
 \widehat\iota_{D,\ell}(H)=O_{\ell(H)}.
 \label{p12-83-c20-s3:eq:elevacion-tipificada-w12-w24}
\end{equation}
La composition \(\widehat\iota_{D,\ell}\) évite d'appliquer
\(\iota_D\), dont le domaine est \(\mathcal H(D)\), directement à une hexade
qui vit encore dans \([12]_{\rm HMT}\).

\begin{proposition}[Injectivité du relèvement relatif]
\label{p12-83-c20-s3:prop:inyectividad-elevacion-w12-w24}
L'application \(\widehat\iota_{D,\ell}\) est injective et conserve
l'incidence entre points et hexades après application de \(\ell\).
\end{proposition}

\begin{proof}
L'isomorphisme \(\ell\) est bijectif sur les blocs. D'après le
\cref{p12-83-c20-s3:thm:diseno-residual-dodecada-autonomo}, chaque hexade résiduelle possède une
unique octade de relèvement. Deux hexades HMT ayant la même image auraient la même
hexade résiduelle et, par bijectivité de \(\ell\), seraient égales.
\end{proof}

Il n'existe pas de choix absolu de \((D,\ell)\) antérieur à cette jauge.
L'objet invariant est la classe de la construction sous changement simultané de
dodécade et d'alignement.

\subsection{Loi d'incidence \texorpdfstring{\(4+1\)}{4+1} sur une tétrade}

Pour une tétrade \(B\subset D\), les paramètres dérivés des deux
systèmes de Steiner sont
\begin{equation}
 \lambda_4(W_{12})
 =\frac{\binom81}{\binom21}=4,
 \qquad
 \lambda_4(W_{24})
 =\frac{\binom{20}1}{\binom41}=5.
 \label{p12-83-c20-s3:eq:ley-cuatro-mas-uno-witt}
\end{equation}
Par conséquent, les cinq octades qui contiennent \(B\) se décomposent de manière
unique en
\begin{equation}
 \{O_H:H\in\mathcal H(D),\ B\subset H\}
 \sqcup\{O_B^\perp\}.
 \label{p12-83-c20-s3:eq:descomposicion-cinco-octadas}
\end{equation}
Les quatre premières satisfont \(|O_H\cap D|=6\) ; la cinquième satisfait
\begin{equation}
 O_B^\perp\cap D=B.
 \label{p12-83-c20-s3:eq:octada-transversal-tetrada}
\end{equation}

Pour \(B=B_K\), les quatre relèvements résiduels correspondent aux
pétales
\[
 \{1,3\},
 \qquad
 \{2,11\},
 \qquad
 \{5,7\},
 \qquad
 \{8,12\},
\]
et la cinquième octade est la transversale \(O_{B_K}^\perp\). Cette loi explique
la transition de quatre hexades à cinq octades sans identifier une face
d'un cube à une hexade ni un sommet à une octade.

\subsection{Correspondance des strates et liberté du triplet positif}

La structure à cinq régions du canal \(\pi\) a pour types
\[
 1_\perp+1_{--}+3_{++}.
\]
L'incidence sur \(B_K\) a pour types
\[
 1_{\rm transversal}+1_{H^-}+3_{\rm restantes}
\]
après désignation de l'hexade négative. Il existe donc une
correspondance naturelle entre \emph{strates de rôle} :
\begin{equation}
 1_\perp\longleftrightarrow O_{B_K}^\perp,
 \qquad
 1_{--}\longleftrightarrow
 \widehat\iota_{D,\ell}(H_\alpha^-),
 \qquad
 3_{++}\longleftrightarrow
 \mathcal O_{B_K}^{+,3}.
 \label{p12-83-c20-s3:eq:correspondencia-estratos-cinco-cinco}
\end{equation}
Ici, \(\mathcal O_{B_K}^{+,3}\) est l'ensemble non ordonné des trois
autres octades résiduelles. Une bijection individuelle entre les trois régions
positives et ces trois octades exige de déclarer une jauge supplémentaire de
\(S_3\). La coïncidence du motif \(1+1+3\) ne détermine pas à elle seule cet
ordonnancement.

\subsection{Groupe d'automorphismes du système d'octades}

\begin{theorem}[Groupe de Mathieu \(M_{24}\)]
\label{p12-83-c20-s3:thm:automorfismos-w24-m24}
Le groupe d'automorphismes du système de Steiner des octades est
\begin{equation}
 \boxed{
 \operatorname{Aut}(W_{24})\cong M_{24},
 \qquad
 |M_{24}|=244823040.}
 \label{p12-83-c20-s3:eq:automorfismos-w24-m24}
\end{equation}
\end{theorem}

Le théorème identifie le groupe classique d'automorphismes de la donnée binaire
d'entrée. Il n'affirme pas que \(M_{24}\) soit engendré à partir d'une seule
dodécade, ni que le code ternaire \(C_W\) détermine à lui seul
\(\mathcal G_{24}\).

\subsection{Séparation catégorique par rapport à la branche de réseaux}

La branche construite dans cette unité est
\begin{equation}
 (W_{12}^{\rm HMT};\mathcal G_{24},D,\ell)
 \longrightarrow W_{24}
 \longrightarrow\operatorname{Aut}(W_{24})\cong M_{24}.
 \label{p12-83-c20-s3:eq:rama-binaria-w24-m24}
\end{equation}
La branche de réseaux qui sera construite ensuite est
\begin{equation}
 C_W\longrightarrow N(A_2^{12})
 \xrightarrow{\text{voisin d’indice trois}}\Lambda_{24}.
 \label{p12-83-c20-s3:eq:rama-ternaria-niemeier-leech-anunciada}
\end{equation}
Leurs constructeurs, domaines et codomaines sont distincts. En particulier,
\begin{equation}
 \boxed{
 \text{aucune flèche n’est introduite }
 W_{24}\longrightarrow\Lambda_{24}.}
 \label{p12-83-c20-s3:eq:no-flecha-w24-leech}
\end{equation}

\begin{figure}[htbp]
\centering
\begin{tikzpicture}[
  >=Stealth,
  node distance=13mm and 16mm,
  box/.style={draw,rounded corners=1.5mm,align=center,
    inner xsep=3mm,inner ysep=2.2mm},
  arr/.style={->,thick}
]
\node[box] (cw) {\(C_W\)};
\node[box,right=of cw] (w12) {\(W_{12}\)};
\node[box,right=of w12] (w24) {\(W_{24}\)};
\node[box,right=of w24] (m24) {\(M_{24}\)};
\node[box,below=of cw] (cw2) {\(C_W\)};
\node[box,right=of cw2] (n) {\(N(A_2^{12})\)};
\node[box,right=of n] (leech) {\(\Lambda_{24}\)};
\draw[arr] (cw)--(w12);
\draw[arr] (w12)--node[above,font=\scriptsize]
 {\((\mathcal G_{24},D,\ell)\)} (w24);
\draw[arr] (w24)--(m24);
\draw[arr] (cw2)--node[below,font=\scriptsize]
 {pegado ternario} (n);
\draw[arr] (n)--node[below,font=\scriptsize]
 {\(3\)-voisin} (leech);
\node[below=5mm of w24,font=\footnotesize,align=center]
 {L'alignement vertical n'est pas une flèche.};
\end{tikzpicture}
\caption{Bifurcation typée des prolongations binaire et ternaire.}
\label{p12-83-c20-s3:fig:bifurcacion-w24-leech-autonoma}
\end{figure}

\subsection{Obstructions de la prolongation binaire \(W_{12}\to W_{24}\) : code de Golay, octades, alignement et séparation de la branche de Leech}

\begin{criterion}[Réfutation de la prolongation binaire]
La construction est réfutée si l'un des faits
suivants se produit :
\begin{enumerate}
 \item le spectre des poids de \(\mathcal G_{24}\) contient un poids non
       répertorié dans \eqref{p12-83-c20-s3:eq:pesos-golay-binario-extendido} ;
 \item les supports de poids huit ne forment pas \(S(5,8,24)\) ;
 \item une intersection résiduelle d'une octade avec \(D\) qui contient cinq
       points de \(D\) a un cardinal différent de six ;
 \item une hexade résiduelle admet zéro octade de relèvement ou plus d'une ;
 \item on applique \(\iota_D\) directement à une hexade HMT sans
       l'alignement \(\ell\) ;
 \item la loi d'incidence sur une tétrade n'est pas \(4+1\) ;
 \item on ordonne le triplet positif sans déclarer la jauge \(S_3\) ;
 \item \(\operatorname{Aut}(W_{24})\) n'a pas pour ordre \(244823040\) ;
 \item on affirme que \(W_{24}\) naît de \(W_{12}\) sans la donnée binaire
       \(\mathcal G_{24}\), la dodécade et l'alignement ;
 \item on sérialise les deux branches au moyen d'une flèche
       \(W_{24}\to\Lambda_{24}\).
\end{enumerate}
\end{criterion}


% Unidad U021. Pegado discriminante ternario y Niemeier A2^12.
% Redaccion autonoma desde la autoridad material 1063.

\section{Recollement discriminant ternaire et construction du réseau de Niemeier de type \texorpdfstring{\(A_2^{12}\)}{A₂¹²}}
\label{p12-83-c20-s4:chap:pegado-niemeier-a2-doce}

Cette unité commence la branche de réseaux directement à partir du code ternaire
autodual \(C_W\). Le système binaire \(W_{24}\), la dodécade et le groupe
\(M_{24}\) ne sont pas des entrées de cette construction. La succession typée est
\[
 C_W
 \longrightarrow
 (A_2^*/A_2)^{12}
 \longrightarrow
 N(C_W)
 \cong N(A_2^{12}).
\]

\subsection{Réseau de racines \texorpdfstring{\(A_2\)}{A₂}}

Soit
\[
 A_2=\mathbb Z\alpha_1\oplus\mathbb Z\alpha_2
\]
le réseau de racines dont la matrice de Gram, dans la base ordonnée
\((\alpha_1,\alpha_2)\), est
\begin{equation}
 G_{A_2}=
 \begin{pmatrix}
 2&-1\\
 -1&2
 \end{pmatrix}.
 \label{p12-83-c20-s4:eq:gram-a2-autonomo}
\end{equation}
Pour \(x=x_1\alpha_1+x_2\alpha_2\),
\begin{equation}
 \langle x,x\rangle
 =2(x_1^2-x_1x_2+x_2^2).
 \label{p12-83-c20-s4:eq:norma-a2-autonoma}
\end{equation}
En particulier, \(A_2\) est pair, défini positif et
\[
 \det A_2=3.
\]

On définit
\[
 \rho:=\alpha_1+\alpha_2.
\]
Alors
\begin{equation}
 \rho^2=2,
 \qquad
 \langle\rho,\alpha_1\rangle
 =\langle\rho,\alpha_2\rangle=1.
 \label{p12-83-c20-s4:eq:identidades-rho-a2}
\end{equation}

\subsection{Module et forme discriminants}

Le réseau dual est
\[
 A_2^*
 :=\{x\in A_2\otimes_{\mathbb Z}\mathbb Q:
       \langle x,A_2\rangle\subseteq\mathbb Z\}.
\]
Nous choisissons les représentants
\begin{equation}
 \mu_0=0,
 \qquad
 \mu_1=\frac{2\alpha_1+\alpha_2}{3},
 \qquad
 \mu_2=\frac{\alpha_1+2\alpha_2}{3}.
 \label{p12-83-c20-s4:eq:representantes-discriminantes-a2}
\end{equation}
Leurs accouplements avec la base simple sont
\[
\begin{array}{c|cc}
 &\alpha_1&\alpha_2\\ \hline
 \mu_1&1&0\\
 \mu_2&0&1
\end{array},
\]
de sorte que \(\mu_1,\mu_2\in A_2^*\). En outre,
\[
 3\mu_1=2\alpha_1+\alpha_2\in A_2,
 \qquad
 3\mu_2=\alpha_1+2\alpha_2\in A_2,
\]
tandis que \(\mu_1,\mu_2\notin A_2\). On obtient
\begin{equation}
 A_2^*/A_2
 =\{\overline\mu_0,\overline\mu_1,\overline\mu_2\}
 \cong\mathbb Z/3\mathbb Z
 \cong\mathbb F_3.
 \label{p12-83-c20-s4:eq:discriminante-a2-f3}
\end{equation}

La somme des représentants se réduit explicitement au moyen de
\[
 \mu_1+\mu_1-\mu_2=\alpha_1,
 \qquad
 \mu_1+\mu_2=\rho,
 \qquad
 \mu_2+\mu_2-\mu_1=\alpha_2.
\]
Par conséquent,
\begin{equation}
 \mu_r+\mu_s-\mu_{r+s\,({\rm mod}\,3)}\in A_2
 \qquad(r,s\in\mathbb F_3).
 \label{p12-83-c20-s4:eq:adicion-representantes-discriminantes-a2}
\end{equation}

\begin{proposition}[Forme discriminante de \(A_2\)]
\label{p12-83-c20-s4:prop:forma-discriminante-a2}
Sous \eqref{p12-83-c20-s4:eq:discriminante-a2-f3}, les formes bilinéaire et quadratique
discriminantes sont
\begin{align}
 b(\overline\mu_r,\overline\mu_s)
 &\equiv\frac{2rs}{3}\pmod{\mathbb Z},
 \label{p12-83-c20-s4:eq:bilineal-discriminante-a2}\\
 q(\overline\mu_r)
 &\equiv\frac{2r^2}{3}\pmod{2\mathbb Z}.
 \label{p12-83-c20-s4:eq:cuadratica-discriminante-a2}
\end{align}
\end{proposition}

\begin{proof}
La formule \eqref{p12-83-c20-s4:eq:norma-a2-autonoma} donne
\[
 \mu_1^2=\mu_2^2=\frac23,
 \qquad
 \langle\mu_1,\mu_2\rangle=\frac13.
\]
Pour \(r,s\in\{0,1,2\}\), ces valeurs coïncident, modulo
\(\mathbb Z\), avec \(2rs/3\). La formule quadratique résulte du choix
\(r=s\) en rappelant que le réseau est pair.
\end{proof}

\subsection{\(12\) facteurs et application de recollement}

Soit
\[
 L_0:=A_2^{12}.
\]
Alors
\[
 L_0^*=(A_2^*)^{12},
 \qquad
 L_0^*/L_0\cong\mathbb F_3^{12}.
\]
Pour \(c=(c_1,\ldots,c_{12})\in\mathbb F_3^{12}\), nous définissons
\begin{equation}
 \phi(c):=(\mu_{c_1},\ldots,\mu_{c_{12}})\in L_0^*.
 \label{p12-83-c20-s4:eq:aplicacion-pegado-cw}
\end{equation}
L'extension déterminée par \(C_W\) est
\begin{equation}
 \begin{aligned}
 N(C_W)
 &:=\bigcup_{c\in C_W}(L_0+\phi(c))\\
 &=\{x\in L_0^*:x\bmod L_0\in C_W\}.
 \end{aligned}
 \label{p12-83-c20-s4:eq:definicion-niemeier-cw}
\end{equation}

\begin{lemma}[Clôture additive du recollement]
\label{p12-83-c20-s4:lem:clausura-pegado-cw}
L'ensemble \(N(C_W)\) est un sous-groupe discret de \(L_0^*\) de rang
vingt-quatre.
\end{lemma}

\begin{proof}
Soient \(x=\ell+\phi(c)\) et \(y=m+\phi(d)\), avec \(\ell,m\in L_0\) et
\(c,d\in C_W\). D'après \eqref{p12-83-c20-s4:eq:adicion-representantes-discriminantes-a2},
\[
 \phi(c)+\phi(d)-\phi(c+d)\in L_0.
\]
Puisque \(C_W\) est linéaire, \(c+d\in C_W\), et alors
\(x+y\in L_0+\phi(c+d)\). La même identité appliquée à l'opposé donne
les inverses. Les inclusions
\[
 L_0\subseteq N(C_W)\subseteq L_0^*
\]
prouvent la discrétion et le rang vingt-quatre.
\end{proof}

\subsection{Intégralité induite par l'autodualité}

La forme bilinéaire discriminante de \(L_0\) est la somme orthogonale de douze
copies de \eqref{p12-83-c20-s4:eq:bilineal-discriminante-a2}. Pour
\(c,d\in\mathbb F_3^{12}\),
\begin{equation}
 b(c,d)
 \equiv\frac23\sum_{i=1}^{12}c_id_i
 \pmod{\mathbb Z}.
 \label{p12-83-c20-s4:eq:bilineal-discriminante-a2-doce}
\end{equation}

\begin{lemma}[Intégralité du recollement]
\label{p12-83-c20-s4:lem:integralidad-pegado-cw}
Le réseau \(N(C_W)\) est entier.
\end{lemma}

\begin{proof}
L'autodualité \(C_W=C_W^\perp\) implique
\[
 \sum_{i=1}^{12}c_id_i\equiv0\pmod3
 \qquad(c,d\in C_W).
\]
D'après \eqref{p12-83-c20-s4:eq:bilineal-discriminante-a2-doce}, la forme discriminante
s'annule sur \(C_W\times C_W\). Par conséquent, l'accouplement de deux classes
recollées quelconques est entier.
\end{proof}

\subsection{Parité induite par les poids du code}

Sur une coordonnée non nulle du code, la forme quadratique discriminante
vaut \(2/3\) modulo \(2\mathbb Z\). Par conséquent,
\begin{equation}
 q(c)\equiv\frac23\operatorname{wt}(c)\pmod{2\mathbb Z}.
 \label{p12-83-c20-s4:eq:cuadratica-discriminante-cw}
\end{equation}

\begin{lemma}[Parité du recollement]
\label{p12-83-c20-s4:lem:paridad-pegado-cw}
Le réseau \(N(C_W)\) est pair.
\end{lemma}

\begin{proof}
Les seuls poids de \(C_W\) sont \(0,6,9,12\), et
\[
 \frac23\{0,6,9,12\}=\{0,4,6,8\}\subseteq2\mathbb Z.
\]
Toute classe de recollement est isotrope pour la forme quadratique discriminante,
de sorte que toutes les normes de \(N(C_W)\) sont paires.
\end{proof}

Soient \(g_1,\ldots,g_6\) les lignes de \(G_W=(I_6\mid A_W)\) et
\(\gamma_i=\phi(g_i)\). La multiplication exacte dans \(A_2^*\) donne
\begin{equation}
 (\langle\gamma_i,\gamma_j\rangle)_{i,j=1}^{6}
 =
 \begin{pmatrix}
 4&2&2&2&2&2\\
 2&4&2&2&2&2\\
 2&2&4&2&2&2\\
 2&2&2&4&2&2\\
 2&2&2&2&4&2\\
 2&2&2&2&2&4
 \end{pmatrix}.
 \label{p12-83-c20-s4:eq:gram-generadores-pegado-cw}
\end{equation}
Cette matrice fournit un contrôle direct supplémentaire de l'intégralité et de la
parité pour les six générateurs de recollement.

\subsection{Indice, déterminant et unimodularité}

Les classes latérales de \(L_0\) qui apparaissent dans
\eqref{p12-83-c20-s4:eq:definicion-niemeier-cw} sont indexées par les
\(|C_W|=729=3^6\) éléments du code et sont distinctes. En conséquence,
\begin{equation}
 [N(C_W):L_0]=3^6.
 \label{p12-83-c20-s4:eq:indice-niemeier-sobre-a2-doce}
\end{equation}
Comme
\[
 \det L_0=(\det A_2)^{12}=3^{12},
\]
la formule du déterminant pour une extension finie donne
\begin{equation}
 \det N(C_W)
 =\frac{\det L_0}{[N(C_W):L_0]^2}
 =\frac{3^{12}}{(3^6)^2}=1.
 \label{p12-83-c20-s4:eq:determinante-niemeier-cw}
\end{equation}

\subsection{Système de racines induit par l'incidence de Witt et classification de ses orbites}

Dans chacune des deux classes discriminantes non nulles de
\(A_2^*/A_2\), la norme minimale est
\[
 \mu_1^2=\mu_2^2=\frac23.
\]
Si \(c\in C_W\setminus\{0\}\), alors
\(\operatorname{wt}(c)\geq6\). Par orthogonalité des douze facteurs,
tout vecteur de \(L_0+\phi(c)\) a une norme au moins égale à
\begin{equation}
 6\cdot\frac23=4.
 \label{p12-83-c20-s4:eq:cota-clases-no-triviales-niemeier}
\end{equation}
Par conséquent, aucune classe de recollement non triviale ne contient de vecteurs de norme
deux. Les racines de \(N(C_W)\) sont exactement les racines déjà présentes dans
\(L_0=A_2^{12}\).

\begin{theorem}[Réalisation du Niemeier de type \(A_2^{12}\)]
\label{p12-83-c20-s4:thm:niemeier-a2-doce-autonomo}
Le réseau \(N(C_W)\) est défini positif, pair et unimodulaire, et son système
de racines est \(A_2^{12}\). En conséquence,
\begin{equation}
 \boxed{N(C_W)\cong N(A_2^{12}).}
 \label{p12-83-c20-s4:eq:identificacion-niemeier-a2-doce}
\end{equation}
\end{theorem}

\begin{proof}
La positivité provient de \(L_0^*\). La
\cref{p12-83-c20-s4:lem:integralidad-pegado-cw,p12-83-c20-s4:lem:paridad-pegado-cw} démontre
l'intégralité et la parité ; \eqref{p12-83-c20-s4:eq:determinante-niemeier-cw} démontre
l'unimodularité. La borne \eqref{p12-83-c20-s4:eq:cota-clases-no-triviales-niemeier}
identifie le système de racines. La classification de Niemeier associe de
manière unique à ce système de racines le réseau indiqué.
\end{proof}

\subsection{Obstructions du recollement ternaire \(C_W\to N(A_2^{12})\) : parité, unimodularité et système de racines}

\begin{criterion}[Réfutation du recollement ternaire]
La construction est réfutée si l'un des faits
suivants se produit :
\begin{enumerate}
 \item l'un des représentants \(\mu_0,\mu_1,\mu_2\) n'appartient pas à
       \(A_2^*\), deux de leurs classes coïncident, ou
       \eqref{p12-83-c20-s4:eq:adicion-representantes-discriminantes-a2} échoue ;
 \item la forme discriminante n'est pas celle donnée dans
       \cref{p12-83-c20-s4:prop:forma-discriminante-a2} ;
 \item \(N(C_W)\) n'est pas stable par addition ou passage à l'inverse ;
 \item il existe \(c,d\in C_W\) avec
       \(\frac23\sum_i c_id_i\notin\mathbb Z\) ;
 \item un mot de \(C_W\) produit
       \(\frac23\operatorname{wt}(c)\notin2\mathbb Z\) ;
 \item l'indice n'est pas \(3^6\) ou le déterminant n'est pas un ;
 \item une classe de recollement non triviale contient un vecteur de norme deux ;
 \item on utilise \(W_{24}\), une dodécade ou une donnée binaire dans l'une
       quelconque des opérations de recollement précédentes.
\end{enumerate}
\end{criterion}


% Unidad U022. Origen incidencial, vecino ternario y Leech.
% Redaccion autonoma desde la autoridad material 1063.

\section{Origine d'incidence, problème radial typé et construction du voisin unimodulaire sans racines}
\label{p12-83-c20-s5:chap:origen-incidencial-vecino-leech}

Le réseau \(N(A_2^{12})\) est déjà construit. Cette unité détermine un
point de base à partir d'un cadre ordonné d'incidences, résout un
problème variationnel dans un domaine déclaré, vérifie les trois conditions
du voisin d'indice trois et démontre que le résultat est le réseau de
Leech. Le vecteur du voisin est noté \(v_{\mathcal F_*}\) : il dépend du
cadre d'incidence et n'intervient pas dans la retenue numérique de \(\alpha\).

\subsection{Cadre ordonné d'incidences}

Considérons les quatre hexades
\begin{align}
 H^-&=\{4,5,6,7,9,10\},
 \label{p12-83-c20-s5:eq:hexada-marco-negativa}\\
 H_e&=\{1,3,4,6,9,10\},
 \label{p12-83-c20-s5:eq:hexada-marco-e}\\
 H_\varphi&=\{1,2,3,4,7,9\},
 \label{p12-83-c20-s5:eq:hexada-marco-phi}\\
 H_{\rm APP}&=\{2,3,5,10,11,12\}.
 \label{p12-83-c20-s5:eq:hexada-marco-app}
\end{align}
La donnée est le cadre ordonné
\begin{equation}
 \mathcal F:=(H^-,H_e,H_\varphi,H_{\rm APP}).
 \label{p12-83-c20-s5:eq:marco-incidencial-ordenado}
\end{equation}
On définit
\begin{equation}
 o(\mathcal F)
 :=(H_e\cap H_\varphi)\setminus(H^-\cup H_{\rm APP}).
 \label{p12-83-c20-s5:eq:operador-origen-incidencial}
\end{equation}

\begin{theorem}[Reconstruction équivariante de l'origine]
\label{p12-83-c20-s5:thm:origen-incidencial-equivariante}
Pour le cadre canonique \(\mathcal F_*\),
\begin{equation}
 \boxed{o(\mathcal F_*)=\{1\}.}
 \label{p12-83-c20-s5:eq:origen-incidencial-singleton}
\end{equation}
En outre, pour toute permutation simultanée \(g\in M_{12}\),
\begin{equation}
 o(g\mathcal F)=g\,o(\mathcal F).
 \label{p12-83-c20-s5:eq:equivarianza-origen-incidencial}
\end{equation}
\end{theorem}

\begin{proof}
On a
\[
 H_e\cap H_\varphi=\{1,3,4,9\}
\]
et
\[
 (H^-\cup H_{\rm APP})\cap\{1,3,4,9\}=\{3,4,9\}.
\]
La différence est \(\{1\}\). Pour l'équivariance, on applique les
identités
\[
 g(A\cap B)=gA\cap gB,
 \quad
 g(A\cup B)=gA\cup gB,
 \quad
 g(A\setminus B)=gA\setminus gB,
\]
valables pour toute bijection.
\end{proof}

La tétrade \(B_K=\{4,6,9,10\}\) ne détermine pas à elle seule l'origine. Le
singleton s'obtient à partir du cadre ordonné complet ; ce n'est pas une coordonnée
choisie après observation du résultat.

\subsection{Chambre radiale positive du système de racines et sélection du cône dominant}

Soit \(\rho_i\) la copie de \(\rho=\alpha_1+\alpha_2\) dans le
\(i\)-ième facteur de \(A_2^{12}\). Nous considérons le domaine
\begin{equation}
 \mathscr R_+
 :=\left\{
 v(a)=\sum_{i=1}^{12}a_i\rho_i:
 a_i=1+3b_i, b_i\in\mathbb Z_{\geq0}
 \right\}.
 \label{p12-83-c20-s5:eq:camara-radial-positiva}
\end{equation}
Puisque les facteurs sont orthogonaux et \(\rho_i^2=2\), si
\(S=\sum_i b_i\), alors
\begin{align}
 v(a)^2
 &=2\sum_{i=1}^{12}(1+3b_i)^2\notag\\
 &=24+12S+18\sum_{i=1}^{12}b_i^2.
 \label{p12-83-c20-s5:eq:norma-radial-expandida}
\end{align}
En réduisant modulo dix-huit,
\[
 v(a)^2\equiv6+12S\pmod{18},
\]
et donc
\begin{equation}
 v(a)^2\equiv0\pmod{18}
 \quad\Longleftrightarrow\quad
 S\equiv1\pmod3.
 \label{p12-83-c20-s5:eq:congruencia-radial-suma-b}
\end{equation}

\begin{proposition}[Minima dans la chambre radiale déclarée]
\label{p12-83-c20-s5:prop:minimos-camara-radial}
Dans le sous-ensemble de \(\mathscr R_+\) défini par
\(v(a)^2\equiv0\pmod{18}\), la norme minimale est \(54\). Elle est atteinte
exactement en douze vecteurs : pour chaque \(j\in\{1,\ldots,12\}\), le
coefficient \(a_j\) vaut quatre et les onze autres valent un.
\end{proposition}

\begin{proof}
La condition \eqref{p12-83-c20-s5:eq:congruencia-radial-suma-b} et la non-négativité des
\(b_i\) impliquent que la plus petite valeur possible de \(S\) est un. Cela
impose un unique indice \(j\) avec \(b_j=1\) et \(b_i=0\) pour
\(i\neq j\). En substituant dans \eqref{p12-83-c20-s5:eq:norma-radial-expandida},
\[
 v(a)^2=24+12+18=54.
\]
Il y a douze choix de \(j\). La valeur admissible suivante de \(S\) est
quatre et produit une norme strictement supérieure.
\end{proof}

L'origine \(o(\mathcal F_*)=\{1\}\) sélectionne
\begin{equation}
 \boxed{
 v_{\mathcal F_*}
 =4\rho_1+\rho_2+\cdots+\rho_{12},
 \qquad
 v_{\mathcal F_*}^2=54.}
 \label{p12-83-c20-s5:eq:vector-marco-incidencial}
\end{equation}

\begin{remark}[Portée exacte de la minimalité]
La norme \(54\) est minimale dans la chambre radiale positive
\eqref{p12-83-c20-s5:eq:camara-radial-positiva}, non dans toute la classe résiduelle. Le vecteur
\[
 u=(\alpha_1-2\alpha_2,\rho,\ldots,\rho)
\]
satisfait
\[
 \alpha_1-2\alpha_2-\rho=-3\alpha_2,
\]
de sorte que
\[
 u\equiv(\rho,\ldots,\rho)
 \equiv v_{\mathcal F_*}\pmod{3A_2^{12}}.
\]
Cependant,
\[
 (\alpha_1-2\alpha_2)^2
 =2(1+2+4)=14,
 \qquad
 u^2=14+11\cdot2=36.
\]
Ce témoin empêche d'étendre l'affirmation de minimalité hors du domaine
de \cref{p12-83-c20-s5:prop:minimos-camara-radial}.
\end{remark}

\subsection{Caractère et noyau de congruence}

Écrivons
\[
 N:=N(A_2^{12}),
 \qquad
 v:=v_{\mathcal F_*}.
\]
On définit
\begin{equation}
 \chi_v:N\longrightarrow\mathbb Z/3\mathbb Z,
 \qquad
 \chi_v(x)=\langle x,v\rangle\bmod3,
 \label{p12-83-c20-s5:eq:caracter-vecino-leech}
\end{equation}
son noyau
\begin{equation}
 N_0:=\ker\chi_v
 =\{x\in N:\langle x,v\rangle\equiv0\pmod3\},
 \label{p12-83-c20-s5:eq:nucleo-congruencia-vecino}
\end{equation}
et l'extension
\begin{equation}
 N_v:=N_0+\mathbb Z\frac v3.
 \label{p12-83-c20-s5:eq:definicion-vecino-indice-tres}
\end{equation}

\subsection{Condition I1 : surjectivité et exclusion des racines}

Par construction, \(v\in A_2^{12}\subseteq N\). Les racines d'un facteur
\(A_2\) sont
\[
 \pm\alpha_1,
 \quad
 \pm\alpha_2,
 \quad
 \pm\rho.
\]
Leurs accouplements avec \(\rho\) sont \(\pm1\) ou \(\pm2\). Tous les
coefficients de \(v\) sont congrus à un modulo trois :
\[
 4\equiv1\pmod3,
 \qquad
 1\equiv1\pmod3.
\]
Par conséquent, aucune racine de \(A_2^{12}\) n'a d'accouplement nul avec
\(v\) modulo trois. En particulier,
\[
 \chi_v(\alpha_{1,1})=1,
\]
de sorte que \(\chi_v\) est surjectif et
\begin{equation}
 [N:N_0]=3.
 \label{p12-83-c20-s5:eq:indice-nucleo-vecino}
\end{equation}
En outre, aucune racine de \(N\) n'appartient à \(N_0\).

\subsection{Condition I2 : divisibilité de la norme}

L'identité \eqref{p12-83-c20-s5:eq:vector-marco-incidencial} donne
\begin{equation}
 v^2=54\equiv0\pmod{18},
 \qquad
 \left(\frac v3\right)^2=6.
 \label{p12-83-c20-s5:eq:condicion-i2-vecino}
\end{equation}
En particulier, \(\langle v,v\rangle\equiv0\pmod3\), donc
\(v\in N_0\).

\subsection{Condition I3 : classe duale d'ordre \(3\)}

Si \(x\in N_0\), il existe \(k\in\mathbb Z\) avec
\(\langle x,v\rangle=3k\). Par conséquent,
\[
 \left\langle x,\frac v3\right\rangle=k\in\mathbb Z,
\]
et
\begin{equation}
 \frac v3\in N_0^*.
 \label{p12-83-c20-s5:eq:v-tercio-en-dual-nucleo}
\end{equation}
D'autre part,
\[
 \left\langle\frac v3,\alpha_{1,1}\right\rangle=\frac43.
\]
Puisque \(\alpha_{1,1}\in N\) et que \(N\) est entier, cela démontre
\begin{equation}
 \frac v3\notin N.
 \label{p12-83-c20-s5:eq:v-tercio-no-en-niemeier}
\end{equation}
Enfin, \(3(v/3)=v\in N_0\). La classe de \(v/3\) dans
\(N_0^*/N_0\) est non triviale et a pour ordre exactement trois. En
conséquence,
\begin{equation}
 [N_v:N_0]=3.
 \label{p12-83-c20-s5:eq:indice-extension-vecino}
\end{equation}

\subsection{Parité, intégralité et déterminant du voisin}

\begin{proposition}[Structure unimodulaire du voisin]
\label{p12-83-c20-s5:prop:estructura-unimodular-vecino}
Le réseau \(N_v\) est entier, pair et unimodulaire.
\end{proposition}

\begin{proof}
Puisque \(N\) est unimodulaire et \([N:N_0]=3\),
\[
 \det N_0=[N:N_0]^2\det N=9.
\]
Comme \([N_v:N_0]=3\),
\[
 \det N_v=\frac{\det N_0}{[N_v:N_0]^2}=1.
\]

Soit \(y=x+m(v/3)\), avec \(x\in N_0\), et écrivons
\(\langle x,v\rangle=3k\). Alors
\begin{align*}
 y^2
 &=x^2+2m\left\langle x,\frac v3\right\rangle
       +m^2\left(\frac v3\right)^2\\
 &=x^2+2mk+6m^2.
\end{align*}
Les trois termes forment un entier pair. La forme polarisée prouve
l'intégralité, et le déterminant un prouve l'unimodularité.
\end{proof}

\subsection{\(2\) classes globales et \(6\) classes locales}

Globalement, \eqref{p12-83-c20-s5:eq:definicion-vecino-indice-tres} n'ajoute que deux classes
latérales non triviales :
\[
 N_0+\frac v3,
 \qquad
 N_0-\frac v3.
\]
Localement, dans chaque facteur \(A_2\), une coordonnée de ces classes appartient
à l'une des six classes
\begin{equation}
 \mathscr C_{j,\pm}
 :=A_2+\mu_j\pm\frac\rho3,
 \qquad j\in\{0,1,2\}.
 \label{p12-83-c20-s5:eq:seis-clases-locales-vecino}
\end{equation}
Dans la première coordonnée, \(4\rho/3\equiv\rho/3\pmod{A_2}\), de sorte que
la même liste vaut pour les douze facteurs.

\begin{table}[htbp]
\centering
\caption{Représentants minimaux des six classes latérales locales.}
\label{p12-83-c20-s5:tab:seis-clases-locales-vecino}
\begin{tabular}{cclc}
\toprule
classe&représentant minimal&\((u,v)\)&
\(2(u^2-uv+v^2)/9\)\\
\midrule
\(\mathscr C_{0,-}\)&\(-\rho/3\)&\((-1,-1)\)&\(2/9\)\\
\(\mathscr C_{0,+}\)&\( \rho/3\)&\((1,1)\)&\(2/9\)\\
\(\mathscr C_{1,-}\)&\( \mu_1-\rho/3\)&\((1,0)\)&\(2/9\)\\
\(\mathscr C_{1,+}\)&\( \mu_1+\rho/3-\rho\)&\((0,-1)\)&\(2/9\)\\
\(\mathscr C_{2,-}\)&\( \mu_2-\rho/3\)&\((0,1)\)&\(2/9\)\\
\(\mathscr C_{2,+}\)&\( \mu_2+\rho/3-\rho\)&\((-1,0)\)&\(2/9\)\\
\bottomrule
\end{tabular}
\end{table}

Ici, \((u,v)\) signifie le représentant
\((u\alpha_1+v\alpha_2)/3\).

\begin{lemma}[Minimum exact des six classes]
\label{p12-83-c20-s5:lem:minimo-seis-clases-locales}
Chaque classe de \eqref{p12-83-c20-s5:eq:seis-clases-locales-vecino} a pour norme minimale
exactement \(2/9\).
\end{lemma}

\begin{proof}
La norme de \((u\alpha_1+v\alpha_2)/3\) est
\[
 \frac29(u^2-uv+v^2).
\]
La forme entière \(u^2-uv+v^2\) est définie positive. Dans une classe résiduelle
non nulle, elle ne prend pas la valeur zéro, donc elle vaut au moins un. Les six
représentants de la \cref{p12-83-c20-s5:tab:seis-clases-locales-vecino} atteignent un.
\end{proof}

\subsection{Critère d'absence de racines dans le réseau de Leech et conséquence pour la distance minimale}

\begin{theorem}[Absence de vecteurs de norme deux]
\label{p12-83-c20-s5:thm:vecino-sin-raices}
Le réseau \(N_v\) ne contient pas de vecteurs de norme deux.
\end{theorem}

\begin{proof}
Toutes les racines de \(N\) sont celles de \(A_2^{12}\). La condition I1 prouve
qu'aucune n'appartient à \(N_0\). La classe \(N_0\) ne contient donc pas
de vecteurs de norme deux.

Dans une nouvelle classe \(N_0\pm v/3\), chacune des douze coordonnées
appartient à l'une des six classes locales. D'après
\cref{p12-83-c20-s5:lem:minimo-seis-clases-locales}, chaque coordonnée contribue pour au moins
\(2/9\). L'orthogonalité donne
\[
 \left\|x\pm\frac v3\right\|^2
 \geq12\cdot\frac29=\frac83>2.
\]
Les nouvelles classes ne contiennent pas non plus de racines.
\end{proof}

\subsection{Identification au réseau de Leech}

\begin{theorem}[Construction du réseau de Leech]
\label{p12-83-c20-s5:thm:identificacion-vecino-leech}
Le réseau
\begin{equation}
 N_v
 =\ker(\langle\,\cdot\,,v_{\mathcal F_*}\rangle\bmod3)
  +\mathbb Z\frac{v_{\mathcal F_*}}3
 \label{p12-83-c20-s5:eq:formula-final-vecino-leech}
\end{equation}
est isométrique au réseau de Leech :
\begin{equation}
 \boxed{N_v\cong\Lambda_{24}.}
 \label{p12-83-c20-s5:eq:identificacion-leech-autonoma}
\end{equation}
\end{theorem}

\begin{proof}
Le réseau \(N_v\) est de rang vingt-quatre et défini positif. La
\cref{p12-83-c20-s5:prop:estructura-unimodular-vecino} démontre qu'il est pair et unimodulaire ;
le \cref{p12-83-c20-s5:thm:vecino-sin-raices} démontre qu'il est dépourvu de racines. Par la
classification de Niemeier, il existe une unique classe d'isométrie de réseaux
définis positifs, pairs et unimodulaires de rang vingt-quatre sans racines :
le réseau de Leech.
\end{proof}

\subsection{Les 243 syndromes du code ternaire poinçonné}

Soit \(G_{11}\) le code ternaire de Golay de paramètres
\([11,6,5]_3\). Son espace de syndromes a pour cardinal
\begin{equation}
 3^{11-6}=3^5=243.
 \label{p12-83-c20-s5:eq:numero-sindromes-golay-ternario}
\end{equation}
Le nombre d'erreurs de poids au plus deux est
\begin{equation}
 1+2\binom{11}{1}+2^2\binom{11}{2}
 =1+22+220=243.
 \label{p12-83-c20-s5:eq:bola-radio-dos-golay-ternario}
\end{equation}

\begin{proposition}[Correspondance parfaite erreur--syndrome]
\label{p12-83-c20-s5:prop:correspondencia-error-sindrome-golay}
Chaque syndrome de \(G_{11}\) possède un unique représentant d'erreur de poids
au plus deux.
\end{proposition}

\begin{proof}
La distance minimale cinq implique que deux erreurs distinctes de poids au
plus deux ne peuvent différer d'un mot du code : leur différence aurait un
poids au plus quatre. L'application erreur--syndrome est donc
injective sur la boule de rayon deux. Les deux ensembles ont pour cardinal
\(243\) d'après \eqref{p12-83-c20-s5:eq:numero-sindromes-golay-ternario} et
\eqref{p12-83-c20-s5:eq:bola-radio-dos-golay-ternario} ; l'application est bijective.
\end{proof}

Ces \(243\) syndromes ne sont pas les \(243\) mots visibles obtenus à partir du
catalogue TPK. Ils partagent un cardinal, mais possèdent des domaines, des équivalences et
des fonctions distincts.

\subsection{Obstructions du voisin ternaire \(N_v\cong\Lambda_{24}\) : classe d'ordre \(3\), absence de racines et séparation des cardinaux \(243\)}

\begin{criterion}[Réfutation du voisin et de l'identification de Leech]
La construction est réfutée si l'un des faits
suivants se produit :
\begin{enumerate}
 \item l'opérateur \(o(\mathcal F)\) ne produit pas \(\{1\}\) ou n'est pas
       équivariant ;
 \item il existe dans la chambre radiale un vecteur admissible de norme inférieure à
       \(54\), ou le nombre de minimiseurs n'est pas douze ;
 \item on affirme une minimalité globale malgré le témoin de norme \(36\) ;
 \item \(\chi_v\) n'est pas surjectif ou une racine appartient à \(N_0\) ;
 \item \(v^2\not\equiv0\pmod{18}\) o \((v/3)^2\neq6\);
 \item \(v/3\notin N_0^*\), \(v/3\in N\), ou sa classe n'est pas d'ordre trois ;
 \item le déterminant de \(N_v\) n'est pas un ou le voisin cesse d'être pair ;
 \item l'une des six classes locales a un minimum inférieur à \(2/9\) ;
 \item l'une des deux nouvelles classes globales contient une racine ;
 \item on confond les deux classes globales avec les six classes locales ;
 \item on utilise \(W_{24}\) comme antécédent du recollement ou du voisin ;
 \item les \(243\) syndromes sont identifiés aux \(243\) mots
       visibles au motif qu'ils partagent le même cardinal.
\end{enumerate}
\end{criterion}


\section{Holonomie finie, configuration de Schläfli et symétrie \texorpdfstring{\(E_6\)}{E6}}
\label{p12-83-c20-s6:chap:holonomia-schlafli-e6-nuevo}
\index{holonomie finie}
\index{Schläfli}
\index{E6@\(E_6\)}

La branche de Witt n'épuise pas les structures exceptionnelles accessibles à partir d'APP. La branche parallèle de ce chapitre part des deux polarisations mixtes somme--produit et produit--somme, avec une orientation de plaquette et une section centrale déclarées. La courbure calculée détermine la forme alternée ; celle-ci définit l'extension centrale. Une extension n'est pas inférée de la seule cardinalité 27.

\subsection{Courbures mixtes d’APP}

Sur \(T_9=(\mathbb Z/9\mathbb Z)^2\), avec
\(S(x,y)=x+y\) et \(P(x,y)=xy\), soient
\[
 \Delta_xf=f(x+1,y)-f(x,y),\qquad
 \Delta_yf=f(x,y+1)-f(x,y).
\]
Les connexions mixtes sont
\[
 A^{SP}=(\Delta_xS,\Delta_yP)=(1,x),\qquad
 A^{PS}=(\Delta_xP,\Delta_yS)=(y,1).
\]
Leur courbure discrète est
\[
 F_A=\Delta_xA_y-\Delta_yA_x.
\]

\begin{proposition}[Courbures orientées opposées]
\label{p12-83-c20-s6:prop:curvaturas-app-e6-nuevo}
\[
 F_{A^{SP}}=1,\qquad F_{A^{PS}}=-1\pmod9.
\]
La circulation sur un rectangle orienté de côtés \(r,s\) est,
respectivement, \(rs\) et \(-rs\).
\end{proposition}

\begin{proof}
\(\Delta_xx-\Delta_y1=1\) et
\(\Delta_x1-\Delta_yy=-1\). Un rectangle contient \(rs\) plaquettes ; lors de la
sommation, les arêtes intérieures s’annulent et il reste la circulation du bord.
\end{proof}

Pour des déplacements \(u=(a,b)\), \(v=(c,d)\), la classe alternée est
\begin{equation}
 \sigma(u,v)=ad-bc\pmod9.
 \label{p12-83-c20-s6:eq:forma-alternante-e6-nueva}
\end{equation}
L’orientation fixe le signe de \(\sigma\). La section centrale employée
ci-dessous appartient à la jauge de cette réalisation ; la modifier par un cobord
produit une extension isomorphe, non un nouvel invariant.

\subsection{Extension centrale et réduction ternaire}

Nous définissons le groupe de Heisenberg nonadique
\begin{equation}
 \mathcal H_9=\{(a,b,z):a,b,z\in\mathbb Z/9\mathbb Z\},
 \quad
 (a,b,z)(c,d,w)=(a+c,b+d,z+w+bc).
 \label{p12-83-c20-s6:eq:heisenberg-9-nuevo}
\end{equation}
Le commutateur de deux déplacements horizontaux est
\[
 [(a,b,0),(c,d,0)]=(0,0,bc-ad),
\]
qui enregistre \(-\sigma(u,v)\). En réduisant chaque coordonnée modulo trois, on
obtient
\begin{equation}
 H_3=\{(a,b,z):a,b,z\in\mathbb F_3\},\qquad |H_3|=27,
 \label{p12-83-c20-s6:eq:heisenberg-3-nuevo}
\end{equation}
avec la même loi. Le noyau de \(\mathcal H_9\to H_3\) est \((3\mathbb Z/9\mathbb Z)^3\), de cardinal 27.

\subsection{Représentation de Weyl--Heisenberg et phase centrale}

L’extension nonadique admet une réalisation unitaire explicite. Soit
\[
 \zeta_9:=\exp(2\pi i/9),
 \qquad
 \mathscr H_9^{\rm Sch}:=\mathbb C^9,
\]
de base \(\{|n\rangle:n\in\mathbb Z/9\mathbb Z\}\). Nous définissons
\begin{equation}
 X|n\rangle:=|n+1\rangle,
 \qquad
 Z|n\rangle:=\zeta_9^n|n\rangle.
 \label{p12-83-c20-s6:eq:weyl-desplazamiento-fase-u029}
\end{equation}
On a
\[
 ZX=\zeta_9XZ
\]
et, par multiplication,
\begin{equation}
 (X^aZ^b)(X^cZ^d)
 =\zeta_9^{bc}X^{a+c}Z^{b+d}.
 \label{p12-83-c20-s6:eq:cociclo-weyl-nonadico-u029}
\end{equation}

\begin{theorem}[Représentation fidèle de Weyl--Heisenberg]
\label{p12-83-c20-s6:thm:representacion-weyl-heisenberg-u029}
L'application
\begin{equation}
 \rho_{\rm WH}:\mathcal H_9\longrightarrow
 U(\mathscr H_9^{\rm Sch}),
 \qquad
 \rho_{\rm WH}(a,b,z):=\zeta_9^zX^aZ^b,
 \label{p12-83-c20-s6:eq:representacion-weyl-heisenberg-u029}
\end{equation}
est une représentation unitaire fidèle. Pour
\(u=(a,b)\) et \(v=(c,d)\), son commutateur satisfait
\begin{equation}
 [\rho_{\rm WH}(u,0),\rho_{\rm WH}(v,0)]
 =\zeta_9^{\,bc-ad}I
 =\zeta_9^{-\sigma(u,v)}I.
 \label{p12-83-c20-s6:eq:conmutador-weyl-holonomia-u029}
\end{equation}
\end{theorem}

\begin{proof}
L’identité \eqref{p12-83-c20-s6:eq:cociclo-weyl-nonadico-u029} reproduit le terme
\(bc\) de \eqref{p12-83-c20-s6:eq:heisenberg-9-nuevo} ; d’où la propriété
d’homomorphisme. Si \(\zeta_9^zX^aZ^b=I\), l’action sur la base impose
successivement \(a=0\), \(b=0\) et \(z=0\). Cela prouve la fidélité. Comparer
les produits dans les deux ordres donne
\(\zeta_9^{bc-ad}\), qui est la phase alternée de
\eqref{p12-83-c20-s6:eq:forma-alternante-e6-nueva}.
\end{proof}

La représentation distingue deux observables conjugués. Soient
\[
 Q_{\rm H}:=\operatorname{diag}(0,1,\ldots,8),
 \qquad
 F_{jk}:=\frac13\zeta_9^{jk},
 \qquad
 P_{\rm H}:=F^\dagger Q_{\rm H}F.
\]
Les bases propres de \(Q_{\rm H}\) et \(P_{\rm H}\) sont mutuellement non biaisées, puisque \(|F_{jk}|^2=1/9\). Pour tout vecteur unitaire \(\psi\in\mathscr H_9^{\rm Sch}\),
\begin{align}
 \operatorname{Var}_\psi(Q_{\rm H})
 \operatorname{Var}_\psi(P_{\rm H})
 &\ge
 \frac14\left|
 \langle\psi,[Q_{\rm H},P_{\rm H}]\psi\rangle
 \right|^2,
 \label{p12-83-c20-s6:eq:incertidumbre-weyl-u029}\\
 H_{Q_{\rm H}}(\psi)+H_{P_{\rm H}}(\psi)
 &\ge\log9.
 \label{p12-83-c20-s6:eq:incertidumbre-entropica-weyl-u029}
\end{align}
Ces opérateurs sont adimensionnels. Une carte physique d'unités est une
réalisation ultérieure et n'intervient pas dans l'extension centrale.

\subsection{Action de \(10\) déplacements sur le graphe de Cayley}

Pour une forme linéaire \(\ell(a,b)=pa+qb\), nous posons
\[
 D_\ell(a,b)=\left(a,b,\frac{ab}{2}+\ell(a,b)\right),
 \qquad Z=(0,0,1),
\]
où \(2^{-1}=2\) dans \(\mathbb F_3\). L'ensemble stable par inversion
\begin{equation}
 S_\ell=
 \{D_\ell(v):v\in\mathbb F_3^2\setminus\{0\}\}
 \cup\{Z,Z^{-1}\}
 \label{p12-83-c20-s6:eq:conjunto-conexion-e6}
\end{equation}
possède dix éléments et ne contient pas l'identité.

\begin{theorem}[Indépendance de la jauge linéaire]
\label{p12-83-c20-s6:thm:independencia-calibre-e6}
Les neuf fonctionnelles \(\ell\) produisent des graphes de Cayley isomorphes. Les
neuf ensembles \(S_\ell\) forment une seule orbite sous
\(\operatorname{Aut}(H_3)\).
\end{theorem}

\begin{proof}
L'application
\[
 \phi_\ell(a,b,z)=(a,b,z+\ell(a,b))
\]
est un automorphisme central et satisfait
\(\phi_\ell(S_0)=S_\ell\). Pour écrire la famille complète dans les
coordonnées polarisées de \cref{p12-83-c20-s6:eq:heisenberg-9-nuevo}, soit
\(c(v,w)=b c\) pour \(v=(a,b)\), \(w=(c,d)\). Étant donné
\(M\in\operatorname{GL}(2,3)\), le défaut

\[
 B_M(v,w)=c(Mv,Mw)-\det(M)c(v,w)
\]
est symétrique, car les parties alternées des deux termes coïncident. Nous définissons la correction quadratique
\begin{equation}
 q_M(v)=2B_M(v,v),
 \qquad
 q_M(v+w)-q_M(v)-q_M(w)=B_M(v,w),
 \label{p12-83-c20-s6:eq:correccion-cuadratica-automorfismos-h3}
\end{equation}
où \(2=2^{-1}\) dans \(\mathbb F_3\). La famille complète d'automorphismes s'écrit alors
\begin{equation}
 (v,t)\longmapsto
 \bigl(Mv,\det(M)t+q_M(v)+\ell(v)\bigr),
 \qquad M\in\operatorname{GL}(2,3),\
 \ell\in(\mathbb F_3^2)^*.
 \label{p12-83-c20-s6:eq:automorfismos-h3-completos}
\end{equation}
L’identité polaire de \cref{p12-83-c20-s6:eq:correccion-cuadratica-automorfismos-h3}
prouve directement que le produit central est conservé. La famille possède
\(48\cdot9=432\) éléments. L’énumération exhaustive des
\(\binom{13}{5}=1287\) sous-ensembles stables par inversion de cardinal dix
trouve exactement ces neuf sous-ensembles avec les paramètres dérivés
ci-après.
\end{proof}

Écrivons \(\underline S=\sum_{s\in S}s\) dans
\(\mathbb Z[H_3]\). Le dénombrement des produits ordonnés donne
\begin{equation}
 \underline S^{\,2}
 =10e+\underline S+5(\underline{H_3}-e-\underline S)
 =5\underline{H_3}+5e-4\underline S.
 \label{p12-83-c20-s6:eq:identidad-grupo-schlafli-nueva}
\end{equation}
En effet, l'identité reçoit dix factorisations ; un élément de \(S\) en reçoit une ; chacun des seize autres en reçoit cinq.

\begin{theorem}[Graphe des vingt-sept droites]
\label{p12-83-c20-s6:thm:grafo-27-rectas-nuevo}
La matrice d’adjacence \(A_\ell\) de
\(\Gamma_\ell=\operatorname{Cay}(H_3,S_\ell)\) satisfait
\begin{equation}
 A_\ell^2=5I-4A_\ell+5J.
 \label{p12-83-c20-s6:eq:identidad-schlafli-nueva}
\end{equation}
Par conséquent,
\[
 \Gamma_\ell=\operatorname{srg}(27,10,1,5),
 \qquad
 \operatorname{Spec}(A_\ell)=\{10^1,1^{20},(-5)^6\}.
\]
Il est isomorphe au graphe d'intersection des vingt-sept droites d'une surface cubique lisse.
\end{theorem}

\begin{proof}
La représentation régulière de
\cref{p12-83-c20-s6:eq:identidad-grupo-schlafli-nueva} produit
\cref{p12-83-c20-s6:eq:identidad-schlafli-nueva}. Ses entrées indiquent que tout sommet
possède dix voisins, qu’une paire adjacente possède un voisin commun et qu’une paire non
adjacente en possède cinq. Sur la droite constante, \(A_\ell\) agit par 10 ;
dans son complément, \(J=0\) et
\((A_\ell-I)(A_\ell+5I)=0\). La trace nulle fixe les multiplicités 20 et 6.

Dans le modèle d’une cubique obtenue en éclatant six points de
\(\mathbb P^2\), les 27 classes sont
\[
 E_i,\qquad L_{ij}=H-E_i-E_j,\qquad
 Q_i=2H-\sum_{j\ne i}E_j.
\]
Avec \(H^2=1\), \(E_i^2=-1\), \(H\cdot E_i=0\) et
\(E_i\cdot E_j=0\) pour \(i\ne j\), sa matrice d'intersection possède les
mêmes paramètres. L'identification à la configuration classique utilise le
théorème d'unicité de cette réalisation des 27 droites
\cite{Manivel2005} ; elle n'est pas attribuée à un certificat inexistant.
\end{proof}

L’assertion d’isomorphie peut être rendue entièrement effective par
une bijection calculée. L’étiquetage n’est pas absolu : une autre section linéaire
ou un automorphisme de la surface produit un autre tableau de la même classe.
\begin{longtable}{@{}c@{\,}c@{\,}c@{\qquad}l@{}}
\caption{Bijection explicite entre \(H_3\) et les vingt-sept classes de
droites.}
\label{p12-83-c20-s6:tab:h3-27-rectas-u029}\\
\toprule
\(a\)&\(b\)&\(z\)&classe de droite\\
\midrule
\endfirsthead
\multicolumn{4}{@{}l}{\small\itshape Tableau \thetable\ (suite)}\\
\toprule
\(a\)&\(b\)&\(z\)&classe de droite\\
\midrule
\endhead
0&0&0&\(E_1\)\\
0&0&1&\(L_{12}\)\\
0&0&2&\(Q_2\)\\
0&1&0&\(L_{13}\)\\
0&1&1&\(Q_1\)\\
0&1&2&\(E_3\)\\
0&2&0&\(Q_3\)\\
0&2&1&\(E_2\)\\
0&2&2&\(L_{23}\)\\
1&0&0&\(L_{14}\)\\
1&0&1&\(L_{35}\)\\
1&0&2&\(L_{26}\)\\
1&1&0&\(L_{36}\)\\
1&1&1&\(L_{24}\)\\
1&1&2&\(L_{15}\)\\
1&2&0&\(L_{25}\)\\
1&2&1&\(L_{16}\)\\
1&2&2&\(L_{34}\)\\
2&0&0&\(Q_4\)\\
2&0&1&\(L_{46}\)\\
2&0&2&\(E_6\)\\
2&1&0&\(E_5\)\\
2&1&1&\(Q_6\)\\
2&1&2&\(L_{56}\)\\
2&2&0&\(L_{45}\)\\
2&2&1&\(E_4\)\\
2&2&2&\(Q_5\)\\
\bottomrule
\end{longtable}

Comme contrôle local, les dix voisins de \(e=(0,0,0)\) sont les éléments de
\(S_0\). Le tableau les transforme en les dix classes qui rencontrent \(E_1\) :
les cinq \(L_{1j}\) et les cinq \(Q_j\), avec \(2\le j\le6\). La
vérification complète compare les 729 entrées des deux matrices
d’adjacence.

Le complément \(B_\ell=J-I-A_\ell\) vérifie
\begin{equation}
 B_\ell^2=8I+2B_\ell+8J,
 \qquad
 \operatorname{srg}(27,16,10,8),
 \qquad
 \operatorname{Spec}(B_\ell)=\{16^1,4^6,(-2)^{20}\}.
 \label{p12-83-c20-s6:eq:complemento-schlafli}
\end{equation}
C'est le graphe de Schläfli dans la convention d'adjacence complémentaire ; \(A_\ell\) est le graphe d'intersection de paramètres \((27,10,1,5)\).

Chaque arête appartient à un unique triangle, puisque \(\lambda=1\). Le nombre
de triangles est
\[
 \frac{27\cdot10}{6}=45.
\]
Le groupe complet des automorphismes est d’ordre
\[
 27\cdot1920=51840
\]
et, par son action sur la configuration de droites, il est identifié à \(W(E_6)\), et pas seulement par égalité des ordres \cite{Manivel2005}.

\subsection{Le réseau de racines qui réalise \texorpdfstring{\(E_6\)}{E6}}

Pour expliciter l’objet sur lequel agit ce groupe, soit
\begin{equation}
 \operatorname{Pic}(X)
 =\mathbb ZH\oplus\bigoplus_{i=1}^6\mathbb ZE_i,
 \qquad
 H^2=1,\qquad E_i^2=-1,\qquad H\cdot E_i=E_i\cdot E_j=0\ (i\ne j),
 \label{p12-83-c20-s6:eq:picard-cubica-e6}
\end{equation}
le réseau de Picard de l’éclatement de six points généraux de
\(\mathbb P^2\). Sa classe canonique est
\[
 K_X=-3H+E_1+\cdots+E_6.
\]
Le complément orthogonal
\begin{equation}
 L_{E_6}=K_X^\perp\subset\operatorname{Pic}(X),
 \qquad (x,y)_{E_6}:=-x\cdot y,
 \label{p12-83-c20-s6:eq:reticulo-raices-e6}
\end{equation}
est défini positif de rang six. Une base de racines simples est
\begin{equation}
 \begin{aligned}
 \alpha_1&=E_1-E_2,& \alpha_2&=E_2-E_3,&
 \alpha_3&=E_3-E_4,\\
 \alpha_4&=E_4-E_5,& \alpha_5&=E_5-E_6,&
 \alpha_6&=H-E_1-E_2-E_3.
 \end{aligned}
 \label{p12-83-c20-s6:eq:raices-simples-e6}
\end{equation}
Chaque \(\alpha_i\) est orthogonal à \(K_X\), est de norme deux et leurs produits
sont \(-1\) exactement pour les paires adjacentes du diagramme
\(E_6\) : la chaîne \(1-2-3-4-5\), avec le nœud \(6\) relié au nœud \(3\).
Par conséquent, leur matrice de Gram est la matrice de Cartan de \(E_6\).

\begin{proposition}[Les soixante-douze racines et leurs réflexions]
\label{p12-83-c20-s6:prop:setenta-dos-raices-e6}
Les racines de \(L_{E_6}\) sont
\begin{equation}
 \begin{aligned}
 &E_i-E_j &&(i\ne j),\\
 &\pm(H-E_i-E_j-E_k) &&(1\le i<j<k\le6),\\
 &\pm(2H-E_1-\cdots-E_6),
 \end{aligned}
 \label{p12-83-c20-s6:eq:enumeracion-raices-e6}
\end{equation}
au total \(30+40+2=72\). Pour chaque racine \(\alpha\),
\begin{equation}
 s_\alpha(x)=x+(x\cdot\alpha)\alpha
 \label{p12-83-c20-s6:eq:reflexion-raiz-e6}
\end{equation}
préserve l’intersection et \(K_X\). Les six réflexions simples engendrent
\(W(E_6)\) et permutent fidèlement les 27 classes
\(E_i,L_{ij},Q_i\) de \cref{p12-83-c20-s6:thm:grafo-27-rectas-nuevo}.
\end{proposition}

\begin{proof}
Substituer chaque classe de \cref{p12-83-c20-s6:eq:enumeracion-raices-e6} dans \cref{p12-83-c20-s6:eq:picard-cubica-e6} donne \(K_X\cdot\alpha=0\) et \(\alpha^2=-2\). Le décompte des classes est celui indiqué. La formule de réflexion est la réflexion orthogonale pour la forme \(-(\cdot)\) ; elle préserve \(K_X\) car \(K_X\cdot\alpha=0\). Le système simple de \cref{p12-83-c20-s6:eq:raices-simples-e6} engendre toutes les racines et son groupe de Coxeter est \(W(E_6)\). Son action sur les 27 classes de courbes exceptionnelles conserve l'intersection, de sorte qu'elle coïncide avec l'action sur le graphe déjà construit \cite{Manivel2005}.
\end{proof}

\subsection{Stratification des 729 mots}

Nous écrivons un mot de six trits sous la forme
\[
 w=(\ell_1,h_1,\ell_2,h_2,\ell_3,h_3),
\]
et nous définissons deux éléments de \(H_3\),
\[
 L(w)=(\ell_1,\ell_2,\ell_3),\qquad
 H(w)=(h_1,h_2,h_3),
\]
en utilisant la loi centrale dans ses produits. Selon
\(L(w)^{-1}H(w)\), les 729 mots se divisent en
\begin{align*}
 \mathcal D_\ell&=\{w:L(w)^{-1}H(w)=e\},\\
 \mathcal I_\ell&=\{w:L(w)^{-1}H(w)\in S_\ell\},\\
 \mathcal S_\ell&=\{w:L(w)^{-1}H(w)\notin\{e\}\cup S_\ell\}.
\end{align*}

\begin{theorem}[Partition Hensel--Schläfli]
\label{p12-83-c20-s6:thm:particion-hensel-schlafli-nueva}
\begin{equation}
 729=27+270+432,
 \label{p12-83-c20-s6:eq:particion-729-e6-nueva}
\end{equation}
correspondant à la diagonale, à l’incidence et au gauchissement.
\end{theorem}

\begin{proof}
Pour chacune des 27 valeurs de \(L(w)\), il y a une unique valeur de
\(H(w)\) dont la différence est l’identité, dix dont la différence est dans \(S_\ell\) et
seize autres. On multiplie \(27\) par \(1,10,16\).
\end{proof}

Le cardinal \(432\) de la strate gauchie coïncide numériquement avec la
multiplicité de la région transversale de \(\pi\) dans
\cref{eq:descomposiciones-cinco-regiones-pi}. Cela n’identifie pas les deux objets : ici,
on compte des mots classés par \(L(w)^{-1}H(w)\), tandis que là,
on compte des préimages d’un registre \(U_6\). Sans opérateur d’entrelacement explicite et
équivariant entre ces domaines, \(432=432\) n’est qu’un contrôle cardinal.

\subsection{Nécessité de la phase centrale}

L’énumération de tous les sous-ensembles stables par inversion de cardinal dix
dans les cinq groupes d’ordre 27 donne
\[
\begin{array}{c|c|c}
\text{groupe}&\text{abélien}&
\#\operatorname{srg}(27,10,1,5)\\ \hline
C_{27}&\text{oui}&0\\
C_9\times C_3&\text{oui}&0\\
C_3^3&\text{oui}&0\\
C_9\rtimes C_3&\text{non}&9\\
H_3&\text{non}&9.
\end{array}
\]
Ainsi, conserver 27 étiquettes mais abélianiser la composition détruit la configuration. Ni la cardinalité ni le degré ne suffisent.

Si \(\mathcal T\) est l’ensemble des 45 triangles, le cubique
\[
 \mathcal C_\Gamma(x_1,\ldots,x_{27})
 =\sum_{\{i,j,k\}\in\mathcal T}x_ix_jx_k
\]
récupère l'adjacence, car chaque arête appartient à un triangle unique. Son stabilisateur par permutations est exactement
\[
 \operatorname{Aut}(\Gamma_\ell)\cong W(E_6).
\]

\begin{figure}[htbp]
\centering
\[
 \text{courbures APP }\pm1
 \longrightarrow\sigma
 \longrightarrow\mathcal H_9
 \longrightarrow H_3
 \longrightarrow\Gamma_{27}
 \longrightarrow\mathcal T_{45}
 \longrightarrow W(E_6).
\]
\caption{Chaîne causale de la branche d’holonomie mixte. Elle est parallèle à
\(C_W\to W_{12}\to M_{12}\) et n’en est pas un maillon.}
\label{p12-83-c20-s6:fig:cadena-holonomia-e6-u029}
\end{figure}

\begin{criterion}[Critères de réfutation de la branche \(E_6\)]
\label{p12-83-c20-s6:crit:refutacion-rama-e6-u029}
La branche est réfutée si les courbures ne sont pas opposées, si le cocycle
central est effacé, si la correction quadratique
\eqref{p12-83-c20-s6:eq:correccion-cuadratica-automorfismos-h3} ne conserve pas le produit,
si \eqref{p12-83-c20-s6:eq:identidad-grupo-schlafli-nueva} ne produit pas les coefficients
\((10,1,5)\), si un groupe abélien produit la même configuration, si la
partition n’a pas pour somme 729, si le tableau
\ref{p12-83-c20-s6:tab:h3-27-rectas-u029} n’entrelace pas les matrices d’adjacence ou si
\(W(E_6)\) est identifié uniquement par son ordre et non par son action sur les
vingt-sept classes et les 72 racines.
\end{criterion}


\section{Couplage orienté APP–Fock et rigidité transversale de \(12\) fibres}
\label{p12-83-c20-s7:chap:acoplamiento-app-fock}
\index{APP--Fock!couplage orienté}
\index{Fock!degré deux}
\index{orientation de feuillet}

Le déterminant cyclotomique ordinaire ne distingue pas \(C\) de \(C^{-1}\).
L'information somme--produit est conservée en introduisant explicitement la
parité de feuillet et une droite d'orientation. Le couplage résultant
produit l'indice \(54\) à partir de l'agrégat APP. Une fois fixées les quatre
conditions du couplage, \(54\), \(J\) et \(196884\) ne sont pas des entrées de
l'évaluation. Cette indépendance évaluative n'affirme ni l'unicité universelle
de l'ansatz ni l'indépendance historique de sa conception.

\subsection{Cadre triangulaire dans chaque orbite}

Soit
\[
M=\mathbb Z[\mathbb Z/9\mathbb Z]
\]
et soit \(M_{\mathrm{bal}}^{C_3}\) le sous-module des éléments invariants
\(\nu\) dont les six coefficients unitaires coïncident :
\[
\nu_u=c(\nu)
\qquad
\bigl(u\in(\mathbb Z/9\mathbb Z)^\times\bigr).
\]
Les coefficients de \(e_0,e_3,e_6\) demeurent libres.

Pour une orbite libre
\(\mathcal O=\{u,4u,7u\}\), posons
\[
b_v=e_v-e_{4v}\in A(\mathcal O).
\]
Si \(b_v^\flat\) est le covecteur défini par la forme de \(A_2\), les trois
racines du cadre triangulaire satisfont
\[
\boxed{
\sum_{v\in\mathcal O}b_v\otimes b_v^\flat
=3I_{A(\mathcal O)}.}
\]
En conséquence,
\[
\sum_{v\in\mathcal O}\nu_v\,b_v\otimes b_v^\flat
=3c(\nu)I.
\]
Le facteur trois provient de la résolution de l'identité par le cadre de
racines.

\subsection{Couplage orienté \(\mathcal H_{\Sigma,\Pi}^{(m)}\) entre le bilan \(C_3\), la parité de feuillet et le degré \(2\) de l'espace de Fock}

Pour \(m\in2\mathbb N_{>0}\), soit
\[
\begin{aligned}
\mathcal V_m&=\mathcal V_{\Sigma,m}\oplus\mathcal V_{\Pi,m},\\
\mathcal V_{\Sigma,m}&=(A_2\otimes\mathbb C)^{\oplus m/2},\qquad
\mathcal V_{\Pi,m}=(A_2\otimes\mathbb C)^{\oplus m/2},\\
g_m&=C^{\oplus m/2}\oplus(C^{-1})^{\oplus m/2},\\
\mathsf E_m&=I_{\mathcal V_{\Sigma,m}}
\oplus(-I_{\mathcal V_{\Pi,m}}).
\end{aligned}
\]
La spécialisation \(m=12\), avec les étiquettes de paire, de feuillet et d'orbite,
est \(\mathcal V_{12}=W_{\mathrm{APP}}\otimes\mathbb C\). Cette notation
évite de confondre l'espace de représentation avec le système de Witt
\(W_{12}=S(5,6,12)\). Le degré deux de l'espace de
Fock et l'action induite sont
\[
\mathcal F_2(\mathcal V_m)
=\mathcal V_m[2]\oplus\operatorname{Sym}^2(\mathcal V_m[1]),
\]
\[
\Gamma_2(g_m)
=g_m|_{\mathcal V_m[2]}
\oplus\operatorname{Sym}^2(g_m|_{\mathcal V_m[1]}).
\]

Soit \(o_{\Sigma\Pi}\) la droite d'orientation de signe
\(\Pi-\Sigma\), et soit \(M_{\mathrm{bal}}^{C_3,-}\) la copie orientée dans
laquelle l'échange de feuillets agit par \(\nu\mapsto-\nu\). Nous définissons
\[
\boxed{
\begin{aligned}
\mathcal H_{\Sigma,\Pi}^{(m)}&:
M_{\mathrm{bal}}^{C_3,-}\longrightarrow
\operatorname{End}(\mathcal F_2(\mathcal V_m))\otimes o_{\Sigma\Pi},\\
\mathcal H_{\Sigma,\Pi}^{(m)}(\nu)
&=3c(\nu)
\left(0_{\mathcal V_m[2]}\oplus\operatorname{Sym}^2\mathsf E_m\right)
\otimes o_{\Sigma\Pi}.
\end{aligned}}
\]
La transformation est linéaire en \(c\), est portée par
\(\operatorname{Sym}^2(\mathcal V_m[1])\), utilise la parité de feuillet et adopte la
normalisation du cadre triangulaire.

Soit \(\sigma\) l'échange de feuillets. Alors
\[
\nu\longmapsto-\nu,\qquad
o_{\Sigma\Pi}\longmapsto-o_{\Sigma\Pi},\qquad
\mathsf E_m\longmapsto-\mathsf E_m,
\]
et
\[
\operatorname{Sym}^2(-\mathsf E_m)
=\operatorname{Sym}^2(\mathsf E_m).
\]
Par conséquent,
\[
\mathcal H_{\Sigma,\Pi}^{(m)}(\sigma\nu)
=\sigma\mathcal H_{\Sigma,\Pi}^{(m)}(\nu).
\]
Lors de l'inversion de l'orientation, le covecteur dual change également de signe ; ainsi
\((-o_{\Sigma\Pi}^{\vee})(-o_{\Sigma\Pi})=1\), et l'évaluation orientée
demeure normalisée.

Une fois
\(o_{\Sigma\Pi}^{\vee}(o_{\Sigma\Pi})=1\) fixé, la trace orientée est
\[
\operatorname{Tr}_{o,\mathcal F_2}
(T\otimes o_{\Sigma\Pi})
=\operatorname{Tr}_{\mathcal F_2}(T).
\]

\begin{theorem}[Indice orienté APP--Fock]
\label{p12-83-c20-s7:thm:indice-orientado-app-fock}
Supposons que \(m\) fibres \(A_2\), avec \(m\) pair, se répartissent également
entre les orientations \(C\) et \(C^{-1}\). Alors
\[
\boxed{
\operatorname{Tr}_{o,\mathcal F_2}
\left(
\mathcal H_{\Sigma,\Pi}^{(m)}(\nu)\Gamma_2(g_m)
\right)
=-\frac{3m\,c(\nu)}2.}
\]
Pour l'agrégat produit par APP,
\[
\delta=\nu_\Pi-\nu_\Sigma
=12e_0+3e_3+3e_6
-3\sum_{u\in(\mathbb Z/9\mathbb Z)^\times}e_u,
\]
on a \(c(\delta)=-3\) et
\[
m=3\ \text{paires}\cdot2\ \text{feuillets}\cdot2\ \text{orbites}=12.
\]
Par conséquent,
\[
\boxed{
\operatorname{Tr}_{o,\mathcal F_2}
\left(
\mathcal H_{\Sigma,\Pi}^{(12)}(\delta)\Gamma_2(g_{12})
\right)=54.}
\]
\end{theorem}

\begin{proof}
Soit \(A=\mathsf E_mg_m\). Les contributions des deux feuillets s'annulent :
\(\operatorname{tr}A=0\). Puisque \(\mathsf E_m\) commute avec \(g_m\) et
\(\mathsf E_m^2=I\),
\[
\operatorname{tr}(A^2)=\operatorname{tr}(g_m^2)=-m.
\]
L’identité
\[
\operatorname{tr}(\operatorname{Sym}^2A)
=\frac{(\operatorname{tr}A)^2+\operatorname{tr}(A^2)}2
\]
donne \(-m/2\) ; le cadre apporte \(3c(\nu)\).
\end{proof}

\begin{corollary}[Minimalité du degré bosonique détecteur]
\label{p12-83-c20-s7:cor:minimalidad-grado-dos-app-fock-u030}
Dans la classe de couplages déclarée, le degré un ne détecte pas l'agrégat
orienté et le degré deux est le premier degré bosonique positif qui le
détecte :
\[
 \operatorname{tr}(\mathsf E_mg_m)=0,
 \qquad
 \operatorname{tr}\!\left(
 \operatorname{Sym}^2(\mathsf E_mg_m)
 \right)=-\frac m2\neq0.
\]
\end{corollary}

\begin{proof}
La première égalité est l'annulation des deux feuillets. La seconde a été
obtenue dans la preuve du \cref{p12-83-c20-s7:thm:indice-orientado-app-fock}. Puisque le degré
zéro ne contient que l'unité et que le degré un donne une trace nulle, le degré deux est
le premier à fournir une réponse non nulle.
\end{proof}

\subsection{Rigidité transversale du nombre de fibres}

Pour un entier \(m\ge1\), nous définissons la série formelle
\begin{equation}
 T_m(q):=
 q^{-1}\prod_{n\ge1}(1+q^n+q^{2n})^{-m}
 \in q^{-1}\mathbb Z[[q]]
 \label{p12-83-c20-s7:eq:familia-fock-m-fibras-u030}
\end{equation}
et, dans la somme orthogonale \(A_2^m\), le vecteur radial
\begin{equation}
 v_m:=4\rho_1+\rho_2+\cdots+\rho_m,
 \qquad
 \|\rho_j\|^2=2.
 \label{p12-83-c20-s7:eq:familia-radial-m-fibras-u030}
\end{equation}
Nous construisons deux fonctions entières sans utiliser préalablement \(m=12\) :
\[
 F(m):=[q]T_m(q),
 \qquad
 R(m):=\|v_m\|^2.
\]

\begin{proposition}[Coefficient de Fock et norme radiale]
\label{p12-83-c20-s7:prop:dos-funciones-rigidez-doce-u030}
Pour tout \(m\ge1\),
\begin{equation}
 F(m)=\frac{m(m-3)}2,
 \qquad
 R(m)=2(m+15).
 \label{p12-83-c20-s7:eq:funciones-rigidez-doce-u030}
\end{equation}
\end{proposition}

\begin{proof}
Jusqu'à l'ordre deux,
\[
 (1+q+q^2)^{-m}
 =1-mq+\frac{m(m-1)}2q^2+O(q^3),
\]
\[
 (1+q^2+q^4)^{-m}=1-mq^2+O(q^4).
\]
Les facteurs avec \(n\ge3\) ne contribuent pas au coefficient de \(q^2\).
Puisque le facteur extérieur de \(T_m\) est \(q^{-1}\),
\[
 F(m)=\frac{m(m-1)}2-m=\frac{m(m-3)}2.
\]
Par orthogonalité des \(m\) copies de \(A_2\),
\[
 R(m)=4^2\|\rho_1\|^2+
 \sum_{j=2}^{m}\|\rho_j\|^2
 =32+2(m-1)=2(m+15).
\]
\end{proof}

\begin{theorem}[Rigidité cyclotomique--radiale]
\label{p12-83-c20-s7:thm:rigidez-transversal-doce-u030}
L'égalité \(F(m)=R(m)\) équivaut à
\begin{equation}
 (m-12)(m+5)=0.
 \label{p12-83-c20-s7:eq:ecuacion-rigidez-doce-u030}
\end{equation}
Dans le domaine \(m\in\mathbb Z_{>0}\), l'unique solution est
\[
 \boxed{m=12,\qquad F(12)=R(12)=54.}
\]
\end{theorem}

\begin{proof}
Égaler les deux expressions de
\eqref{p12-83-c20-s7:eq:funciones-rigidez-doce-u030} produit
\[
 m(m-3)=4(m+15),
\]
c’est-à-dire,
\[
 m^2-7m-60=(m-12)(m+5)=0.
\]
La seconde racine est négative et se situe hors du domaine.
\end{proof}

Cette rigidité appartient à la famille cyclotomique--radiale déclarée. Elle n'utilise pas la
valeur \(54\) pour choisir \(m\) : elle construit d'abord \(F(m)\) et \(R(m)\), puis
résout leur égalité.

\subsection{Factorisation du couplage APP–Fock : hypothèses, invariants et conditions nécessaires de validité et obstructions}

Le couplage se factorise par
\[
M_{\mathrm{bal}}^{C_3,-}/\ker c
\cong\mathbb Z\otimes o_{\Sigma\Pi},
\]
et l'application induite est injective puisque
\(\operatorname{Sym}^2\mathsf E_m\neq0\).
La formule est sélectionnée par les quatre conditions déclarées :
linéarité en \(c\), support dans
\(\operatorname{Sym}^2(\mathcal V_m[1])\), parité
de feuillet et normalisation par le cadre triangulaire. On n'affirme pas l'unicité
parmi tous les couplages naturels ni l'indépendance historique de
l'ansatz. Une fois ces conditions fixées, \(54\), \(J\) et \(196884\) ne
sont pas des entrées de l'évaluation.

Trois contrôles montrent que le résultat n'est pas fixé par une
normalisation rétrospective :
\[
(m,c)=(12,-2)\longmapsto36,\qquad
(m,c)=(10,-3)\longmapsto45,
\]
et
\[
\delta_k=\delta+ke_0
\]
laisse invariant \(\mathcal H_{\Sigma,\Pi}^{(12)}\) tandis que le poids numérique
passe de \(54\) à \(54+9k\). En particulier, pour \(k=1\),
\[
\operatorname{Tr}_{o,\mathcal F_2}
\left(
\mathcal H_{\Sigma,\Pi}^{(12)}(\delta+e_0)\Gamma_2(g_{12})
\right)=54,
\qquad
w(\delta+e_0)=63.
\]

\begin{proposition}[Nécessité de la donnée d'orientation]
\label{p12-83-c20-s7:prop:orientacion-app-fock}
Si l'on oublie \(o_{\Sigma\Pi}\), tout scalaire naturel linéaire en
\(\nu_\Pi-\nu_\Sigma\) et ne dépendant que de la classe de la représentation
s'annule.
\end{proposition}

\begin{proof}
\(C\) et \(C^{-1}\) sont conjugués, de sorte que la classe de représentation est
paire sous l'échange de feuillets. L'agrégat
\(\nu_\Pi-\nu_\Sigma\) est impair. Sans la droite d'orientation, une
transformation simultanément paire et impaire vaut zéro.
\end{proof}

Le théorème établit ainsi le transport exact de l'agrégat vers le degré deux de
Fock. Le défaut local état par état conserve l'obstruction de descente
du \cref{p12-83-c1-s5:cor:obstruccion-isotipica-app-autonoma} ; les deux énoncés appartiennent à
des domaines distincts et sont compatibles.


% Unidad U031. Alineamiento APP--Witt, automorfia de orden tres,
% construccion VOA y Moonshine.
